\documentclass[11pt,letterpaper,reqno]{amsart}
% BEGIN INLINE manuscript-style
\makeatletter
\NeedsTeXFormat{LaTeX2e}[2020/10/01]
% Use with \documentclass[11pt,letterpaper,reqno]{amsart}.
% This file controls presentation only. It supplies no mathematical content.

\RequirePackage[T1]{fontenc}
\RequirePackage{lmodern}
\RequirePackage{microtype}
\RequirePackage[letterpaper,margin=1in,headheight=15pt]{geometry}
\RequirePackage{amsmath,amsthm,amssymb,mathtools}
\RequirePackage{mathrsfs}
\RequirePackage{enumitem}
\RequirePackage{perpage}
\RequirePackage{longtable}
\RequirePackage{xcolor}
\definecolor{manuscriptred}{rgb}{0.6,0.1,0.1}
\definecolor{manuscriptblue}{rgb}{0.1,0.3,0.5}
\definecolor{manuscriptmagenta}{HTML}{60445C}
\RequirePackage[breaklinks=true,colorlinks=true,
  linkcolor=manuscriptred,citecolor=manuscriptblue,urlcolor=manuscriptmagenta]{hyperref}
\RequirePackage{cleveref}
\hypersetup{pdfauthor={}}

% Retain the AMS paragraph, title, heading and running-head conventions.
% Equation counters restart each section: 2.i, 2.ii, ... .
% Appendix sections keep their letters: A.i, A.ii, ... .
\numberwithin{equation}{section}
\renewcommand{\theequation}{\thesection.\roman{equation}}
% Prose references include "eq." and omit parentheses around the number.
\renewcommand{\eqref}[1]{\textup{eq.~\ref{#1}}}
\crefname{equation}{eq.}{eqs.}
\Crefname{equation}{eq.}{eqs.}
\creflabelformat{equation}{#2#1#3}
\setcounter{tocdepth}{1}
\renewcommand{\thefigure}{\Roman{figure}}
\renewcommand{\thetable}{\Roman{table}}
\newcommand{\rtneedspace}[1]{\par\begingroup
  \dimen0=#1\relax
  \vskip0pt plus\dimen0\penalty-100
  \vskip0pt plus-\dimen0\vskip\dimen0
  \penalty9999\vskip-\dimen0\vskip0pt\endgroup}
\MakePerPage{footnote}
\renewcommand{\thefootnote}{%
  \ifnum\value{footnote}>9\arabic{footnote}\else\fnsymbol{footnote}\fi}

% Upright statements. Main results are A, B, C; supporting statements
% share a section-local sequence. These end marks indicate typography,
% never proof completeness, validation, or a deferred-proof status.
\theoremstyle{definition}
\newtheorem{ms@theorem}{Theorem}
\renewcommand{\thems@theorem}{\Alph{ms@theorem}}
\newtheorem*{ms@theoremstar}{Theorem}
\newtheorem*{ms@theoremzero}{Theorem 0}
\newtheorem{ms@proposition}{Proposition}[section]
\newtheorem{ms@lemma}[ms@proposition]{Lemma}
\newtheorem{ms@definition}[ms@proposition]{Definition}
\newtheorem{ms@corollary}[ms@proposition]{Corollary}
\newtheorem*{ms@corollarystar}{Corollary}

\newcommand{\ms@majorstart}{%
  \pushQED{\qed}\renewcommand{\qedsymbol}{$\blacksquare\square$}}
\newcommand{\ms@minorstart}{%
  \pushQED{\qed}\renewcommand{\qedsymbol}{$\blacksquare$}}
\newenvironment{theorem}[1][]{\ms@majorstart\ms@theorem[#1]}
  {\popQED\endms@theorem}
\newenvironment{theorem*}[1][]{\ms@majorstart\ms@theoremstar[#1]}
  {\popQED\endms@theoremstar}
\newenvironment{theoremzero}[1][]{\ms@majorstart\ms@theoremzero[#1]}
  {\popQED\endms@theoremzero}
\newenvironment{proposition}[1][]{\ms@minorstart\ms@proposition[#1]}
  {\popQED\endms@proposition}
\newenvironment{lemma}[1][]{\ms@minorstart\ms@lemma[#1]}
  {\popQED\endms@lemma}
\newenvironment{definition}[1][]{\ms@minorstart\ms@definition[#1]}
  {\popQED\endms@definition}
\newenvironment{corollary}[1][]{\ms@majorstart\ms@corollary[#1]}
  {\popQED\endms@corollary}
\newenvironment{corollary*}[1][]{\ms@majorstart\ms@corollarystar[#1]}
  {\popQED\endms@corollarystar}
% The source's lemma/proposition variants are deliberately not carried over:
% choose one marker for each statement type rather than encode proof status.

\theoremstyle{remark}
\newtheorem{remark}[ms@proposition]{Remark}
\newtheorem*{remark*}{Remark}
% The ordinary amsthm proof environment retains its open square.

\crefname{ms@theorem}{theorem}{theorems}
\Crefname{ms@theorem}{Theorem}{Theorems}
\crefname{ms@proposition}{proposition}{propositions}
\Crefname{ms@proposition}{Proposition}{Propositions}
\crefname{ms@lemma}{lemma}{lemmas}
\Crefname{ms@lemma}{Lemma}{Lemmas}
\crefname{ms@definition}{definition}{definitions}
\Crefname{ms@definition}{Definition}{Definitions}
\crefname{ms@corollary}{corollary}{corollaries}
\Crefname{ms@corollary}{Corollary}{Corollaries}
\crefname{remark}{remark}{remarks}
\Crefname{remark}{Remark}{Remarks}

% Bibliography configuration for this manuscript.
\makeatother
% END INLINE manuscript-style


\title[Trajectory and partition stability]{Coupling trajectories and partition stability in lattice Yang--Mills theory}
\date{Working draft, 9 November 2026}
\begin{document}
\raggedbottom
\begin{abstract}
We construct a coupling trajectory and its common vacuum normalization
from admitted small-field formation, coefficient and large-field
closure estimates. Uniform block drift gives a protected terminal
window, arbitrarily fine first-passage cutoffs and the leading
logarithmic ultraviolet law without positivity of each increment.
Frozen-choice analytic estimates give tuned bare couplings under
their stated prefix domains. Component bounds and label counting
control the normalized zero-source history sum; a retained
small-field event supplies the matching partition lower bound.
The local analytic constructions, the one-loop coefficient and
the quantitative history estimates are explicit inputs. No
complex-source stability or continuum limit is asserted here.
\end{abstract}
\maketitle
\tableofcontents

\section{Introduction}

Renormalization produces a new action coefficient before it has
finished restoring the form of the effective measure. To iterate
the construction, that computed coefficient must first be admitted
on the domains required by the remaining operations. A proof that
assumes the next complete effective density while checking this
admission is circular. We give an induction in which the current
history data determine the regular output and its coefficient,
the coefficient extends the numerical frame, and the large-field
operation then restores the next hypothesis. The common vacuum
normalization is completed with the remainder returned by that
last operation.

Two further issues are separate from this induction. A positive
average drift need not make every coupling increment positive;
first passage must allow bounded backward excursions. Also, a
stochastic realization preserves total mass without bounding it.
The partition estimates require quantitative bounds on actual
components, the count of their labels, the retained integrals,
and a lower-bound event with the same vacuum normalization.

The regular-iteration companion \cite{EM} supplies bounds along
prescribed admissible prefixes. Here the one-step coefficient
and closure inputs are used to produce a prefix to which those
bounds apply. The scalar forward-profile estimate in that
companion also has a backward fixed-point application. We keep
this homogeneous profile argument distinct from tuning the full
history-dependent trajectory, whose marked terms have separate
analytic domains.

Ba\l aban's effective-density estimates and stability bounds
provide the conditional renormalization framework
\cite[Theorem~2 and Corollary~3, pp.~263--264]{Bal88} and
\cite[Theorem~1 and (0.1), pp.~355--356]{Bal89}.
The running-coupling hypothesis is retained explicitly in the
latter theorem. Those results do not by themselves supply the
common trajectory or the coefficient identifications required
here. We use actual local formation, geometry, normalization and
history estimates as named hypotheses, and prove the ensuing
induction, arithmetic, analytic tuning and partition deductions.

The partition bounds are uniform in cutoff and compatible finite
volume at a fixed sufficiently small terminal coupling. The trajectory
and tuning estimates state their dependence on the inverse-coupling floor. Complex-source estimates, transfer to infinite-volume
Gibbs states and limiting field constructions require further
results. The present arguments do not establish those conclusions.

\section{Finite measures and the common vacuum normalization}
\label{rt:sec-normalization}

Fix an odd integer $L\ge3$, a length $\ell>0$, and a finite
cutoff index $K'\in\mathbb Z_{\ge0}$.
The physical torus is $T_\ell=(\mathbb R/\ell\mathbb Z)^4$.
At level $j$ its lattice $\Lambda_j$ has mesh $L^{j-K'}$, for
$0\le j\le K'$. We require $\ell L^{K'-j}$ to be a positive integer and
require all partitions used by the construction to be compatible
with these periods. Write
\[
 a_{K'}=L^{-K'},\qquad
 N_j=|\Lambda_j|=\ell^4L^{4(K'-j)}.
\]
The geometric admission conditions below may require a larger
period than this divisibility condition alone.

A configuration assigns $U_b\in\operatorname{SU}(2)$ to each
positive bond, with $U_{\bar b}=U_b^{-1}$ on the reverse bond.
Every unoriented plaquette is counted once, with one fixed
orientation. Set
\begin{equation}\label{rt:eq-wilson-measure}
 A(U)=\sum_p\left(1-\tfrac12\operatorname{Tr}U_p\right),
 \qquad d\mu_{0,\beta}(U)=e^{-\beta A(U)}\prod_b dU_b,
 \qquad Z_{K',\ell,\beta}=\mu_{0,\beta}(1),
\end{equation}
where $dU_b$ is Haar probability measure and $\beta>0$.
The trace of an $\operatorname{SU}(2)$ matrix is real. On complexified
link charts we use this same trace polynomial as the holomorphic
extension of the action.
The finite partition function is positive: the integrand is a
strictly positive continuous function on a compact product.
The initial presentation has this Wilson action, one empty history,
and no effective regular, remainder or boundary families. Its
membership in the admitted initial class also requires the
geometric and bare-boundary inputs stated with that class.

At a later level the state includes the current bond configuration,
the retained fields, and the discrete history records. We use the
Borel sigma field on each finite-dimensional field space and the
disjoint-union sigma field for the records. A stochastic kernel
is a measurable probability kernel on these spaces. For a finite
positive measure $\mu$ and a kernel $Q$, the measure $\mu Q$ is
defined by $(\mu Q)(F)=\int Q(F)(u)\,d\mu(u)$ for bounded Borel
$F$. Each fixed real operation used here is supplied as such a
kernel or as a composition of the operations in the next lemma.
An analytic expansion of its density may have summands of either
sign; that does not change the underlying positive measure.

\begin{lemma}[Mass and retained variables]\label{rt:lem-mass}
Suppose a finite positive measure is transformed by finitely many
of the following operations: a Borel test with finitely or countably many outcomes that
partition the entire state space; a draw from a stochastic kernel; or a
measurable map, including a map which forgets variables. Sum over
all test outcomes. The total mass is unchanged.

More specifically, let $m$ be a finite positive measure on
$\mathsf A\times\mathsf B$ and let $\kappa$ be a stochastic
kernel from $\mathsf A$ to $\mathsf C$. Adjoin a draw from
$\kappa$ and then retain only the $\mathsf C$ coordinate. The
result is
\begin{equation}\label{rt:eq-standing-kernel}
       \bigl[m(\kappa\otimes1)\bigr]_{\mathsf C}
                          =m_{\mathsf A}\kappa.
\end{equation}
These statements hold for a fixed signed or complex measure if
all displayed integrals are absolutely integrable.
\end{lemma}
\begin{proof}
For a test with indicators $\chi_r$, the sum of the masses is
$\int\sum_r\chi_r\,d\mu=\mu(1)$ because $\sum_r\chi_r=1$.
For a draw, $Q1=1$ gives $(\mu Q)(1)=\mu(1)$. For a map $f$,
$f_*\mu(1)=\mu(1)$. Composition proves the first assertion.
To prove \eqref{rt:eq-standing-kernel}, test both sides on a
bounded Borel function $F$ on $\mathsf C$. The left side is
\[
 \int_{\mathsf A\times\mathsf B}
       \left(\int_{\mathsf C}F(c)\,\kappa(a,dc)\right)m(da,db).
\]
The inner integral depends only on $a$, so integration against the
marginal $m_{\mathsf A}$ gives the right side. Fubini's theorem
applies because the measure is finite and $F$ is bounded. Absolute
integrability gives the same argument for signed or complex measures.
\end{proof}

For actual Borel block averages $M_j$ between the successive
link spaces, write $T_j=(M_j)_*$. Their bare pushforwards satisfy
\[
 T_{n-1}\cdots T_0=(M_{n-1}\circ\cdots\circ M_0)_*.
\]
Indeed both sides integrate a bounded Borel function against its
composition with $M_{n-1}\circ\cdots\circ M_0$. For $n=0$ both
are the identity. This concerns the bare averaging maps; the
full real operations also keep the stipulated history variables.

Undefined local means, logarithms or minimizers are assigned a
failed-test outcome, so they do not leave an unassigned part of
the configuration space. A gauge or frame draw precedes every
test that depends on that draw, and its value is retained in the
history until it is forgotten by a specified map. Replacing such
a draw by multiplication by a gauge-fixing weight is valid only
under the separate density identity that justifies that replacement.
In particular, the mass passing a subsequent frame-dependent test
must be averaged with the frame actually used by that test.

For a fixed real run let $Q_j$ be its complete operation on the
extended states at step $j\to j+1$, and put
$\mu_{j+1}=\mu_jQ_j$. Its stochastic realization is part of the
actual local construction. If reference laws are fixed in advance,
the same $Q_j$ acts linearly on a family of input measures. An
operation whose reference law is recomputed from each density
cannot be declared linear for this reason. At zero source,
\cref{rt:lem-mass} gives $\mu_j(1)=Z_{K',\ell,\beta}$.
This identity supplies no bound uniform in the cutoff after a vacuum
normalization: every multiple $c\nu$ of a probability measure
also satisfies the same mass-preservation identity, for any $c>0$.

\begin{lemma}[Bounded insertions and positive history kernels]
\label{rt:lem-history-insertion}
Let \(\mu\) be a finite positive measure with positive mass on
a standard Borel bare state space \(X\). Let \(K_h\) be positive measurable
kernels from \(X\) to standard Borel retained state spaces \(Y_h\),
indexed by a finite or countable set of records, with
\(\sum_h K_h(u,Y_h)=1\) for every \(u\). For a bounded Borel
function \(F:X\to\mathbb C\), put
\[
 \tau_h(B)=\int K_h(u,B)\,d\mu(u),\qquad
 \tau_h^F(B)=\int F(u)K_h(u,B)\,d\mu(u).
\]
Then \(|\tau_h^F|\le\|F\|_\infty\tau_h\). The density
\(F_h=d\tau_h^F/d\tau_h\) is a version of the conditional
expectation of \(F\) given the record and retained state, under
the probability law on the disjoint union of the spaces \(X\times Y_h\)
whose record-\(h\) part is \(\mu(du)K_h(u,dy)/\mu(1)\). In particular,
\begin{equation}\label{rt:eq-history-insertion}
  \frac{\int F\,d\mu}{\mu(1)}
    =\frac{\sum_h\int F_h\,d\tau_h}{\sum_h\tau_h(1)},
   \qquad |F_h|\le\|F\|_\infty\quad \tau_h\text{-a.e.}
\end{equation}
If \(f_h\) is Borel and \(F(u)=f_h(y)\) almost everywhere in the joint law of record
\(h\), then \(F_h=f_h\) almost everywhere on \(Y_h\).
If a measurable symmetry \(r\) with measurable inverse preserves \(\mu\), permutes the records
and transports their spaces and kernels by
\(K_{rh}(ru,rB)=K_h(u,B)\), then \(r_*\tau_h=\tau_{rh}\)
and \(r_*\tau_h^F=\tau_{rh}^{F\circ r^{-1}}\).
\end{lemma}
\begin{proof}
For each measurable partition of \(B\), the sum of the moduli
of the inserted integrals is at most
\(\int |F(u)|K_h(u,B)\,d\mu(u)\). The variation bound follows.
The Radon--Nikodym theorem gives \(F_h\) and its essential
bound. Testing the defining identity against functions of
\((h,y)\) identifies the conditional expectation. It also
gives the asserted version when \(F\) is already determined
by the retained state. The sums are absolutely convergent
by the variation bound and \(\sum_h K_h(u,Y_h)=1\). Summing
the inserted and uninserted masses gives
\eqref{rt:eq-history-insertion}. For the covariance assertions,
change \(u\) to \(ru\) in the measures of \(rB\) for record \(rh\).
\end{proof}

For the Wilson measure one may take \(F(u)\) to be a bounded
gauge-invariant function of the iterated bare average. Its
conditional value is the corresponding function of the
retained fields when those fields determine it. The complete
history operation may replace other fields, so the insertion
formula uses the joint law of the original field and the
retained state.

We now fix the numerical subtraction used throughout the paper.
The new regular family at level $j+1$ is
$\mathcal E_{j+1}(X,P,z)$, and the new remainder family is
$\mathcal R_{j+1}(X,P)$. Their carrier $X$ is a specified finite
union of localization cells, $z$ is a pointing index, and
$P=(V,J)$ is the admitted field-current pair. The trivial pair is
$P_\circ=(\mathbf1,0)$. Values at this pair use the actual
localization and its fixed logarithm branches. The sums below
are finite or absolutely convergent in the supplied rooted norms.
They range over the complete periodic reference families, with
the same convention at every history.

Let $\mathfrak n_j$ be the field-independent numerical factor
assigned to the actual step. Define
\begin{equation}\label{rt:eq-vacuum-recursion}
 \begin{aligned}
 E_{(0)}&=0,\qquad E_{(j+1)}=E_{(j)}+\mathsf e_j,\\
 \mathsf e_j&=\mathfrak n_j+
   \sum_X\left\{\sum_z\mathcal E_{j+1}(X,P_\circ,z)
                            +\mathcal R_{j+1}(X,P_\circ)\right\}.
 \end{aligned}
\end{equation}
The number $E_{(j)}$ may depend on the run, cutoff and torus.
It is common to all histories on that fixed torus. Missing local
vacuum contributions in a particular history remain in its
normalized activity; they are not subtracted a second time.
The new remainder in \eqref{rt:eq-vacuum-recursion} is the one
returned by the completed large-field operation. Thus this
recursion is completed after coupling admission and closure,
and is not an input to the calculation of that coupling.

At step $j$, the positive number $x_j$ is the action weight,
and the positive number $\theta_j$ is its frozen reference weight.
For $\tau_j=x_j/\theta_j$, one assigned Gaussian factor in a
resolved $\operatorname{SU}(2)$ fibre is
$-\tfrac32|\mathfrak X_j|\log\tau_j$, where
$\mathfrak X_j$ is its set of free non-pivot bonds.
Each such bond has three real coordinates. The logarithm is
based at $\tau_j=1$ on the admitted connected relative disc.
Haar, gauge and reference-normalizer constants are assigned by
the same actual formula. We use the Gaussian family normalized
at the trivial datum, as in \cite[\S4]{EM}; its removed reference
values are included once in $\mathfrak n_j$. A presentation
retaining the original Gaussian vacuum values is equivalent
only after the opposite change to $\mathfrak n_j$, as proved
below. The gauge-compensation factor
in an interacting cut remains part of that integral; its
unrestricted Gaussian integral cannot replace it there.

\rtneedspace{14\baselineskip}
\begin{lemma}[Common normalization and redistribution of constants]
\label{rt:lem-vacuum-normalization}
Let the actual finite measures and the real numbers in
\eqref{rt:eq-vacuum-recursion} be supplied as above. Put
\[
 \nu_j=e^{-E_{(j)}}\mu_j,\qquad
 \widehat Z_j=e^{-E_{(j)}}Z_{K',\ell,\beta}.
\]
Then
\begin{equation}\label{rt:eq-normalized-operation}
 \nu_{j+1}=e^{-\mathsf e_j}\nu_jQ_j,
 \qquad \nu_j(1)=\widehat Z_j.
\end{equation}
Replacing a local family by that family minus a field-independent
number, and adding the sum of the removed numbers to
$\mathfrak n_j$, leaves $\mathsf e_j$, $E_{(j)}$ and these
normalized measures unchanged. The same assertion holds if the
numbers are moved in the opposite direction.
\end{lemma}
\begin{proof}
Since $Q_j$ is fixed for this run, multiplication of its input
measure by the scalar $e^{-E_{(j)}}$ commutes with its action.
The identity $E_{(j+1)}=E_{(j)}+\mathsf e_j$ therefore gives
the first equation. Taking total mass and using
\cref{rt:lem-mass} gives the second. In
\eqref{rt:eq-vacuum-recursion}, subtracting constants from the
vacuum values cancels their sum added to $\mathfrak n_j$.
Absolute convergence permits that redistribution. Induction on
$j$ leaves every accumulated number unchanged.
\end{proof}

Redistribution is made in the complete density. If a history
omits a carrier, its residual numerical factor changes by the
corresponding omitted constant, so the physical measure stays fixed.

The coupling subtraction has a different role. The coefficient
$b_{j+1}$ is read from the newly formed complete regular family,
using the normalized classical insertion whose coefficient is
one. Subtracting $b_{j+1}$ times that insertion changes the
classical exponent from $-x_j A$ to
$-(x_j-b_{j+1})A$. Its vacuum value is zero, since $A(\mathbf1)=0$.
The new young remainder is retained in its designated class;
its vacuum value contributes to \eqref{rt:eq-vacuum-recursion}
after closure. One does not add its quadratic coefficient as a
second coupling subtraction. Older admitted remainders already
enter the computed full regular coefficient through the actual
regular step.

The counterterm bookkeeping also applies to a spatially varying
record of completed steps. Put
$A(w,U)=\sum_p w(p)(1-\tfrac12\operatorname{Tr}U_p)$ for a plaquette
weight $w$. If the actual history assigns a counterterm $b_i\phi_i(p)$
at plaquette $p$, set $w_0(p)=\beta$ and
$w_i(p)=w_{i-1}(p)-b_i\phi_i(p)$. At a fixed read of the field $U$,
linearity in the weights gives
\[
       -A(w_k,U)-\sum_{i=1}^k b_i A(\phi_i,U)=-\beta A(U).
\]
The profile $\phi_i$ records where step $i$ was performed. On the
all-small-field history it is one for the corresponding complete
counterterm. The displayed identity concerns the actual assignment of
coefficients and does not change the background or replace its tests.

For the real history estimate, the actual profiles on this same
plaquette field are required to satisfy
$1\ge\phi_1(p)\ge\cdots\ge\phi_k(p)\ge0$. If $x_j>0$ and
$b_j=x_{j-1}-x_j$, summation by parts gives
\begin{equation}\label{rt:eq-positive-action-profile}
 w_k(p)=x_0(1-\phi_1(p))
   +\sum_{j=1}^{k-1}x_j(\phi_j(p)-\phi_{j+1}(p))
   +x_k\phi_k(p)\ge0.
\end{equation}
The coefficients are nonnegative and sum to one. This proves
positivity without requiring $b_j\ge0$. For $k=0$ use $w_0=x_0$.
The native geometric condition gives this ordering when each profile
is supported in its action region and equals one on its one-cell
interior, and the next action region lies in that interior. Wherever
the next profile is positive, the preceding one is then one.
On a stock where the actual profiles equal one for $j<t$ and zero
for $j\ge t$, the weight is $x_{t-1}$. That stock identity is a
geometric input about the actual action profiles. All these action
terms use the same field read. An inner retained island uses its
own assigned field and Wilson factor. Every allocation of energy
shares must consume at most the original nonnegative weight on
each plaquette of its own field, leaving a nonnegative residual.

All constants in the later estimates are uniform in the cutoff,
the length of an admitted prefix and the compatible torus,
except where a dependence on the terminal coupling or volume is
displayed. Geometric constants, localization parameters and
analytic reserves are fixed before the final inverse-coupling
floor. A numerical inequality about the resulting increments
does not by itself provide the admitted local maps or their
common domains.

\subsection{The logarithmic cost of a local normalizer}

The numerical factors in a history must be charged on their
actual carriers. One elementary estimate explains the logarithmic
cost of the scalar Haar normalizers used in that assignment.

\begin{lemma}[A truncated Haar normalizer]\label{rt:lem-haar-normalizer}
For $y\ge1$, let $\chi_y:\operatorname{SU}(2)\to[0,1]$ be Borel.
Suppose a fixed $a_*>0$ satisfies
$\chi_y(U)=1$ whenever $\|U-I\|_{\rm op}<a_*y^{-1/2}$.
With Haar probability measure define
\[
 z(y)=\int\chi_y(U)
       e^{-y(1-\frac12\operatorname{Tr}U)}\,dU.
\]
There are constants $0<c\le C<\infty$, depending only on $a_*$,
such that
\[
       c y^{-3/2}\le z(y)\le C y^{-3/2},
       \qquad |\log z(y)|\le C+\tfrac32\log y.\qedhere
\]
\end{lemma}
\begin{proof}
The class angle $\phi\in[0,\pi]$ has Haar density
$(2/\pi)\sin^2\phi\,d\phi$. Also
$1-\tfrac12\operatorname{Tr}U=1-\cos\phi$ and
$\|U-I\|_{\rm op}=2\sin(\phi/2)$.
The inequalities $\sin\phi\le\phi$ and
$1-\cos\phi\ge2\phi^2/\pi^2$ on this interval give
\[
 z(y)\le\frac2\pi\int_0^\infty
                 \phi^2e^{-2y\phi^2/\pi^2}\,d\phi
                      =C y^{-3/2}.
\]
Put $a=\min\{a_*/2,1\}$. For $0\le\phi<a y^{-1/2}$,
the cutoff is one, $\sin\phi\ge2\phi/\pi$, and
$y(1-\cos\phi)\le a^2/2$. Integrating on this interval
gives $z(y)\ge (8a^3/(3\pi^3))e^{-a^2/2}y^{-3/2}$.
Finally $z(y)\le1$ and this lower bound give the logarithmic
estimate.
\end{proof}

For example, a fixed numerical factor $\sigma_0>0$ contributes
$\log\sigma_0-\tfrac32\log x_j$ per assigned free bond,
and a factor $z(x_j)^{-(L^4-1)}$ contributes
$-(L^4-1)\log z(x_j)$ per assigned coarse cell. Each costs
at most $C(1+\log x_j)$ at its own level. Their actual
multiplicities and carriers come from the finite density formula.
When a history omits such a carrier, its cost is charged there;
it is not bounded by the number of final coarse sites without
the geometric and weighted-integral estimates proved below.

The regional formula can be seen explicitly in the
self-calibrated convention \(x_{j-1}=\theta_{j-1}\). At a
step with regions \(\Lambda\subset\Omega\), write \(B_{j-1}^{\rm free}(Y)\)
for the number of incident free fine bonds in \(Y\) after
pivot removal, and \(\#_j(Y)\) for the level-\(j\) site
count. Let \(\sigma_0\) be the Haar density at the identity
in the chosen exponential coordinates, let \(z_{j-1}\) be
the scalar site-gauge normalizer. On $\mathfrak g=\mathfrak{su}(2)$,
with its chosen Euclidean inner product and Lebesgue measure, set
\(\nu_{\mathfrak g}=\log\int_{\mathfrak g}e^{-|Y|^2/2}\,dY\).
In the convention
that keeps the unnormalized Gaussian local vacuum values in
the regular family, the three elementary scalar contributions at this
step are
\begin{equation}\label{rt:eq-native-regional-scalar}
 \begin{aligned}
 &\bigl(\log\sigma_0-\tfrac32\log x_{j-1}\bigr)B_{j-1}^{\rm free}(\Omega)
        -(L^4-1)\#_j(\Omega)\log z_{j-1}\\
 &\hspace{35mm}+\nu_{\mathfrak g}B_{j-1}^{\rm free}(\Lambda).
 \end{aligned}
\end{equation}
The vacuum values of the regular and remainder families on
their actual ranges are added to this number. If a local
Gaussian family is normalized at the trivial pair, its
removed vacuum values are instead added once to the assigned
number, as in \eqref{rt:eq-vacuum-recursion}.

The first term of \eqref{rt:eq-native-regional-scalar}
comes from extracting \(\sigma_0\) from each free-link Haar
density and scaling its three real coordinates by
\(x_{j-1}^{-1/2}\). The second counts the \(L^4-1\)
noncentral site gauges per level-\(j\) block. The last is
the scalar for expressing the fixed reference Gaussian on the
three real coordinates of each counted bond as a probability measure.
This normalization is an identity between reference densities
before the cut is imposed. The cut and every
gauge-compensation factor stay inside the resulting integral.

The last factor normalizes the product standard reference density.
Any other field-independent covariance or pivot number keeps its
prescribed allocation in the raw local Gaussian total or assigned
scalar, with its trivial value counted once. It is not an additional
evaluation of an interacting cut integral.

\begin{proposition}[Regional counts and comparison of prefix scalars]
\label{rt:prop-vacuum-count}
At a step $j-1\to j$, let $Y\subset\Lambda_j$ and identify it with
the union of its fine blocks. Count a positive bond in the regional
normalizer when at least one endpoint is in $Y$, and remove one distinct fine
pivot for each incident coarse bond, using the actual incidence convention. Use $L^4-1$ normalized site
gauges per coarse site in $Y$. Let $d_{j-1}(Y)$ be the number of
remaining incident fine bonds minus the number of these site gauges.
Let $B_j^\partial(Y)$ be the number of coarse bonds with exactly one endpoint
in $Y$. Then
\begin{equation}\label{rt:eq-vacuum-incidence}
\begin{aligned}
 d_{j-1}(\Lambda_j)&=3(L^4-1)N_j,\\
 d_{j-1}(\Lambda_j)-d_{j-1}(Y)
   &=3(L^4-1)|Y^c|-\tfrac12(L^3-1)B_j^\partial(Y).
\end{aligned}
\end{equation}
The second line lies between zero and $3(L^4-1)|Y^c|$.

Suppose the actual normalizers assign the Gaussian scale
$-\tfrac32\log\theta_{j-1}$ to each free bond and the reciprocal
Haar scale $+\tfrac32\log\theta_{j-1}$ to each such gauge. Keep the
remaining scalar prices, including the bounded ratios
$x_{j-1}/\theta_{j-1}$, in the prescribed step constants. The
logarithmic part left after subtracting the full-torus value is
$\tfrac32\log\theta_{j-1}$ times the second line of
\eqref{rt:eq-vacuum-incidence}. In particular it is a nonnegative
cost when $\theta_{j-1}\ge1$.

If the remaining scalar price on a full level is bounded in absolute
value by $C N_j$, and a fixed $C_{\rm exc}\ge0$ satisfies
$\theta_i-\theta_j\ge-C_{\rm exc}$ for $i\le j$, then on
one fixed run, for $0\le k\le n\le K'$,
\begin{equation}\label{rt:eq-vacuum-prefix-comparison}
 E_{(n)}-E_{(k)}
 =-\tfrac92(L^4-1)\sum_{j=k}^{n-1}(\log\theta_j)N_{j+1}
     +O(N_k).
\end{equation}
Consequently replacing the prefix normalization $E_{(k)}$ by
$E_{(n)}$ changes the normalized mass by a factor between
$e^{-C N_k}$ and $e^{C(1+\log\theta_k)N_k}$, for $\theta_j\ge1$.
\end{proposition}
\begin{proof}
There are $4N_j$ positive coarse bonds and $4L^4N_j$ positive
fine bonds. The full free-bond count is $4(L^4-1)N_j$. The full
site-gauge count is $(L^4-1)N_j$, proving the first line.
Count coarse bond endpoints in $Y^c$ to get
$4|Y^c|-B_j^\partial(Y)/2$ bonds with both endpoints outside $Y$. There
are $L^3$ fine crossings per coarse face, so the analogous fine
count is $4L^4|Y^c|-L^3B_j^\partial(Y)/2$. Each omitted pivot corresponds
to a coarse bond with both endpoints outside $Y$. Subtract the
omitted pivot and site-gauge counts from the omitted fine count
to get the second line. The bound $B_j^\partial(Y)\le8|Y^c|$ gives its
nonnegativity for $L\ge3$. Its upper bound follows by dropping the
boundary term. The count applies on the periodic lattice, counting
oriented occurrences if projected edges have the same endpoints.

The signed scale price is $-\tfrac32\log\theta_{j-1}$ times the
difference of free-bond and site-gauge counts. Subtract its value
on the full torus. This gives the sign and magnitude claimed for
the regional omission. Summing the full-torus values in the
once-counted recursion gives the principal term of
\eqref{rt:eq-vacuum-prefix-comparison}. The remaining costs sum
to $O(N_k)$ because $N_{j+1}=L^{-4(j+1-k)}N_k$. Finally
$0\le\log\theta_j\le\log(\theta_k+C_{\rm exc})$ for $j\ge k$.
Exponentiating the comparison with the sign in the normalized
mass proves both factor bounds.
\end{proof}

\section{Joint construction of the weights and presentations}
\label{rt:sec-joint-construction}

Use the finite measures and the vacuum convention of
\cref{rt:sec-normalization}. To distinguish a chosen finite lattice
from the stopping time, write \(K'\) for the finite cutoff and use
\(K\) for first passage. At level \(j\) the period in its own units
is \(P_j=\ell L^{K'-j}\). Prefixes on different tori or cutoffs
are compared through the prescribed common lifted reads.

\subsection{The local inputs}
\label{rt:subsec-local-inputs}

For a real reference frame \(\theta=(\theta_0,\ldots,\theta_k)\),
an admitted presentation through \(k\) records the following data.
It records the weights \(x_j\), their relative domains in the frame,
the complete regular families of both chains and their common
companions, and their coefficients at each created level. It also
records the unpointed remainders with their birth cores, the boundary
terms with their determining data, and the positive history measures
with retained fields, tests, and local exponents. The exponents have
the Wilson weight and the once-counted common scalar of
\cref{rt:eq-vacuum-recursion}. Each regular term in the action is
vacuum subtracted and carries its own classical quadratic
counterterm. Each remainder is vacuum subtracted.

The analytic class includes the declared real gauge covariance.
History kernels are covariant under lattice isometries preserving
all their partitions. Each localized regular or remainder family
has translation covariance under the translations preserving its
own partition. Full lattice-isometry covariance is required of the
unrestricted pointed totals on jointly admitted real background
pairs; it is not imposed on every filed local term. Identifications of complete reads
transport the actual local formulas, including the support tests,
retained variables, measure factors and branches. The assigned
remainder and boundary spaces, together with the birth domains of
standing operations, remain part of the data. A bound on an activity
does not identify its formula, or change the space in which a later
step may read it.

We use three local inputs. They are uniform for the prescribed
constants and for every admitted old prefix, with no hypothesis on a
future trajectory.

\begin{enumerate}
\item The initial Wilson presentation is admitted. At level \(k\),
  the actual regular and boundary formation reads the given
  presentation and the frame through \(k\). It supplies the new
  complete regular output, its absolute-radius unmarked counterpart,
  their common analytic restrictions and companions, and the boundary
  data required for the succeeding operations. The full regular list
  uses the old remainders that have passed their age and domain tests.
  The unmarked list has none. These formations use the same frozen
  reference, real fluctuation cut and physical coordinate choices
  on their common reads. They are available before a new frame entry
  is chosen. In particular the physical creation radii and covariance
  threshold are specified by the old reference \(\theta_k\).
\item The physical coefficient is linear on the common creation
companion. Each finite realization assigns that same companion read
as its counterterm; wrapping terms remain in its finite action bounds.
The unmarked and marked envelopes from the regular
  iteration apply on this prefix \cite[Sections~3--6]{EM}. On the
  self-calibrated choice \(\theta_j=x_j\) for \(j\le k\), the actual
  counterterm has the form
  \begin{equation}\label{rt:eq-trajectory-increment}
     b_{k+1}=b_{k+1}^0+f_{k+1}+q_{k+1},\qquad
     |f_{k+1}|\le\epsilon_{\rm u}(x_k),\quad
     |q_{k+1}|\le\epsilon_{\rm m}(x_k).
  \end{equation}
  The nonnegative envelopes tend to zero at infinity. The actual
  one-loop data are independent of the initial weight, admissible
  torus and unused finer horizon. We require
  \begin{equation}\label{rt:eq-one-loop-input}
     \left|\sum_{l=1}^r(b_l^0-c_0)\right|\le c_2
           \quad(r\ge0),\qquad
     c_0=\frac{11N^2}{12\pi^2}\log L,\qquad N=2 .
  \end{equation}
  Here \(c_2\ge0\). This is the determinant and trace input in the
  Wilson convention \(x_0=\beta=g_0^{-2}\). A Gaussian
  finite-dimensional integration identity alone does not compute
  \(c_0\). For the quantitative estimate below the ordinary regular
  envelopes can be taken as
  \[
     \epsilon_{\rm u}(y)\le C_{\rm u}y^{-1/2}(\log y)^{\mathfrak p_0},
     \qquad
     \epsilon_{\rm m}(y)\le C_{\rm m}y^{-p_R},
     \qquad p_R=(\kappa_0-4)/2 ,
  \]
  on the supplied domains, with fixed \(\mathfrak p_0\ge0\) and
  \(\kappa_0>4\). Dropping the logarithmic improvement in
  the marked estimate only enlarges this envelope.
\item For an admitted presentation and an admitted new weight, the
  actual large-field closure returns the successor presentation
  with the same constants and prolongs the old data. Admission
  requires the real tube, frame inequalities, geometric tests and
  coupling ceilings of the construction. Its conclusion supplies the
  next positive histories, regular and remainder bounds, boundary
  terms, and their covariance and common-read identities. At fixed
  torus and cutoff the actual measure transformation has the mass identity
  of \cref{rt:lem-mass}. The new final remainder and the step scalar
  for \cref{rt:eq-vacuum-recursion} are supplied after this closure.
\end{enumerate}

The common-read input includes the infinite companion and compatibility
under unused finer cutoffs, with the same reference and weight prefix.
For each bare weight above the fixed input floor, the Wilson
presentation with $\theta_0=x_0=\beta$ is admitted on a cofinal family
of sufficiently large tori and fine cutoffs. Given an admitted old
prefix, its regular formation is available on such a family. If its
computed next weight and chosen reference entry satisfy the numerical
admission inequalities, assume there is a cofinal family on which the
entire finite proposed extension has the geometric and analytic read
admissions of the actual closure. This includes the contour, chart,
determining-history and support conditions for the prescribed tests and
outcomes. If two realizations contain the same complete lifted read
with its margins, the prescriptions agree there, including earlier
reads used in the new step. These are local inputs for the initial
entry and for one computed extension at a time. Each realization uses
its own finite measure and scalar constants.

\begin{lemma}[A conditional construction of coherent coefficients]
\label{rt:lem-coherent-supply}
Fix a real frame prefix and its localization parameters. Suppose the
basic local coefficients on compatible tori are holomorphic on the
specified complex read charts, with uniform bounds there. The charts
include the actual determining collars. On each torus take the finite
subgroup of translations preserving the admitted settings, charts,
labels and rooted bounds. The coefficients have a specified sum rule
whose total is covariant under this subgroup.

Suppose a nested cofinal sequence of finite covers eventually covers
each fixed target torus and each bounded lifted read. The covers
supply linear filing maps. A filing assigns lifted labels to their
actual target labels and reads their entries by fixed chart pullbacks.
For each fixed finite part of a filing these pullbacks take compact
subsets of the target chart into compact subsets of the corresponding
source charts. Assume the filings compose under coverings and
restrictions. Each admitted target translation has compatible admitted
lifts on sufficiently large covers, and filing commutes with these
translations. The filed sums obey the actual cover identity for the
prescribed total on each target. Fix $C_0,C_1>0$, $D_{\rm cat}\ge1$ and $0<q<D_{\rm cat}^{-1}$.
On each fixed rooted chart assume there is a fixed catalogue of at most $C_0 D_{\rm cat}^m$ length-$m$ labels,
the same across all covers, with absent entries set to zero and entry bound $C_1q^m$
and filing multiplicity at most $m+1$, uniformly over the
covers. The target charts form a countable collection of finite-dimensional
charts, with countably many
discrete read types, labels and entries.

Then translation averaging preserves the local read, bounds and sum
rule. Along one cofinal subsequence the averaged, filed entries have
holomorphic limits on the common charts. The limits are translation
covariant, compatible under the filing maps, and satisfy the
prescribed sum rules. The compatible restrictions to bounded lifted
reads give a subsequential infinite companion. For representatives in
a box, the comb with a prescribed exterior root and entrance bond
gives a local gauge rule. Transforming the input by a site gauge $u$
transforms its representative by $u$ outside the box and by the
constant $u(o)$ inside, where $o$ is the exterior root.
\end{lemma}
\begin{proof}
On a finite torus let $T$ be the admitted translation subgroup, acting on formulas
by translating the setting, the label and the read coordinates. For
a coefficient family $C$ use $|T|^{-1}\sum_{a\in T}a^*C$. If two
settings have the same complete local read, their translates do also,
so every summand has the same value on the two settings. This proves
locality of the average. Convexity preserves the rooted bounds. A
translation permutes the summands, proving covariance. Summing the
averaged coefficients gives the average of the translated sum rules,
which equals the original covariant total.

For a box representative let $g_U(x)$ be the ordered transport along
the prescribed tree from $o$ to $x$ inside, and put $g_U=1$ outside.
With $U^g(x,y)=g(x)U(x,y)g(y)^{-1}$, each tree bond of $U^{g_U}$
is the identity. The root value and the tree equations determine
$g_U$ uniquely. In the transformed field the tree transport is
$g_{U^u}(x)=u(o)g_U(x)u(x)^{-1}$ inside. The product gauge
$g_{U^u}u$ therefore equals $u(o)g_U$ inside and $u$ outside.
This is the stated transformation law and uses only the actual tree
and its entrance read.

File the averages on the common charts. Uniform boundedness of each
scalar entry gives a normal family on its finite-dimensional chart.
Extract successively on compact exhaustions of the countably many
charts and take a single diagonal subsequence. Each admissible target
translation is realized by the stated lifts for all sufficiently large
covers, so its identity passes to the local uniform limits. The
compact containment for the pullbacks ensures local uniform
convergence of every finite filed sum. For a rooted filing or
sum rule the tail beyond length $s$ is bounded uniformly by
\[
  C_0C_1\sum_{m>s}(m+1)(D_{\rm cat}q)^m,
\]
which tends to zero. Pass first to the limit in each finite partial
sum, then let $s$ increase. Composability of the filing maps before
the limit gives the corresponding identities between limits, and the
supplied cover identity identifies the limiting sum with the
prescribed total on its stated domain.
\end{proof}

The chart margins, compatibility of the filings and cover identity
are inputs about the actual local operators. The limit supplies a
coherent choice of basic coefficients. Exact agreement with an
already fixed finite realization on a nonwrapping read still requires
the stated common-read identification. Identification of the physical
coefficient uses convergence of the actual responses on the common
charts and the moment argument of \cite[Proposition~6.3]{EM}.

\subsection{An order for the constants}

In the local input order, the exponents and fixed separation fractions
precede $L$ and the initial linear cubes. The prescribed decay rates
and their reserves come next, followed by the enlarged localization
cubes and the final cell multiplier $M$. Fix the analytic reserves,
norm caps and cutoff constants before the weak frame triple. The
triple is fixed before the upper coupling ceiling $\gamma$ is chosen.
The bare-disc radius and the room for the admitted background
continuations come after that ceiling, and before the final
inverse-coupling floor and the physical choices of terminal coupling,
compatible torus, bare weight and cutoff. Each local constraint is
checked at its last free parameter after substitution of quantities
determined by earlier choices. For example, a lower bound on a decay
reserve must be met before imposing the larger cube size that pays
for it, and the final coupling floor must allow the already fixed
bare-disc and background rooms. The local analytic and geometric
inequalities used at these stages are the supplied ones.
The reserve for retained integrals is a separate choice. Fix it after
the required loads and their charge constants, and before any coupling
ceiling or inverse-coupling floor used to pay those charges.

The class conditions at a chosen parameter must be finitely many
floors, ceilings, or endpoints of a nonempty window, with earlier
data fixed. A uniform estimate indexed by a later observable
exponent is used only at the finitely many exponents fixed at this
stage.

\rtneedspace{26\baselineskip}
\begin{lemma}[Order of the numerical choices]
\label{rt:lem-constant-order}
Suppose the local inputs admit the preceding finite order of choices.
Fix \(c_0,c_2\) from \cref{rt:eq-one-loop-input}, and choose
\(C_\beta\ge c_0+2c_2+1\). Set
\begin{equation}\label{rt:eq-flow-constants}
 \begin{gathered}
  c_5=2c_2,\qquad \lambda_{\rm w}=c_0/4,\qquad
  B_{\rm w}=C_\beta+1,\\
  \lambda_{\rm s}=c_0/2-\varsigma,\qquad
  B_{\rm s}=C_\beta+\varsigma,\qquad
  0<\varsigma\le\min\{c_0/4,1\}.
 \end{gathered}
\end{equation}
After fixing the weak triple \((\lambda_{\rm w},c_5,B_{\rm w})\),
one can impose all the following numerical conditions on a common
floor \(Y\):
\begin{equation}\label{rt:eq-trajectory-error-floor}
  \epsilon_{\rm u}(y)+\epsilon_{\rm m}(y)\le\min\{c_0/4,1\},
  \qquad \epsilon_{\rm m}(y)\le\varsigma\qquad(y\ge Y).
\end{equation}
It can also impose \(Y>\max\{2c_5,\gamma^{-2}+c_5+1\}\), where
\(\gamma\) is the supplied upper coupling ceiling. For the fixed
cell exponent \(r>0\) it can impose
\[
        \left(\frac{\log(y+2B_{\rm s})}{\log y}\right)^r<L
                  \qquad(y\ge Y).
\]
If the required reserve is
\(\varsigma=C_*\gamma^{\kappa_0-8}\) with \(C_*>0\) fixed before
\(\gamma\), then \(\kappa_0>8\) permits that reserve in
\cref{rt:eq-flow-constants}.
\end{lemma}
\begin{proof}
With earlier data fixed, finitely many floors on a free upward
parameter are combined by their maximum. Finitely many positive
ceilings are combined by their minimum. A window may be chosen
only after its nonemptiness has been checked in the preceding
parameters. Apply this rule through the supplied local choices.
The weak triple is fixed before taking the constants of the marked
bound. If the displayed reserve is used, its positive exponent
allows the ceiling \(\gamma\) to be reduced until \(\varsigma\)
lies in the stated range. Both error envelopes tend to zero with
the fixed earlier constants. This permits \cref{rt:eq-trajectory-error-floor}
and the other finite floors to be met together. The strong triple
satisfies \(\lambda_{\rm s}\ge\lambda_{\rm w}\) and
\(B_{\rm s}\le B_{\rm w}\), so it needs no new estimate on a
different weak class. The logarithmic ratio in the cell
condition tends to one, allowing its floor at the same stage.
\end{proof}

The source uses rounded cells based on a running minimum. Fix the
power \(r>0\) chosen with the geometry, and write
\begin{equation}\label{rt:eq-running-cell}
 R(y)=\min\{L^m:m\in\mathbb Z_{\ge0},\ L^m\ge(\log y)^r\},
 \qquad R_j=R\!\left(\min_{0\le i\le j}x_i\right).
\end{equation}
The real cutoff can be chosen on any cofinal class satisfying the
actual geometric admission conditions. The numerical floor in the
source is \(P_j\ge 2L^{11}M R_j\) at each used level, where \(M\)
is the fixed localization side in level units. If a construction
needs a prospective floor at \(j+1\) before \(x_{j+1}\) is known,
\(P_{j+1}\ge2L^{11}M R_j\) is sufficient for that numerical
condition. The local contour and history-domain tests remain
separate inputs.

\subsection{Admission and closure}

\begin{theorem}[The joint trajectory through first passage]
\label{rt:thm-joint-trajectory}
Assume the three local inputs and the cofinal admission stated
in \cref{rt:subsec-local-inputs} on weak prefixes, with the constants
chosen as in \cref{rt:lem-constant-order}. On an admitted old
prefix the weak frame inequalities are
\[
  \lambda_{\rm w}(j-i)-c_5\le x_i-x_j
                           \le B_{\rm w}(j-i)\qquad(i\le j).
\]
Fix \(h\ge Y\) and \(\beta\ge h+B_{\rm s}\). Start from the
admitted Wilson presentation, put \(x_0=\beta\), and use
\(\theta_j=x_j\) as each new entry is admitted. The prescribed
construction forms the real coefficients and actual presentations
through
\begin{equation}\label{rt:eq-first-passage}
   K=K(\beta,h)=\min\{j\ge0:x_j\le h+B_{\rm s}\}.
\end{equation}
The coefficients and the stopping index are common to all
admitted tori and unused finer horizons that realize this prefix.
The joint vacuum recursion is \cref{rt:eq-vacuum-recursion} with
the same convention in each finite realization.

The stopping index is finite. For all \(0\le i\le j\le K\),
\begin{equation}\label{rt:eq-strong-drift}
 \tfrac34c_0(j-i)-c_5\le x_i-x_j\le B_{\rm s}(j-i),
 \qquad
 \lambda_{\rm s}(j-i)-c_5\le x_i-x_j.
\end{equation}
The terminal value is protected by the preceding window:
\begin{equation}\label{rt:eq-first-passage-bounds}
 h<x_K\le h+B_{\rm s},\qquad
 \frac{\beta-h}{B_{\rm s}}-1\le K
 \le1+\frac{\beta-h-B_{\rm s}+c_5}{\lambda_{\rm s}} .
\end{equation}
In particular, for fixed \(h\) the stopping depths are unbounded
as \(\beta\) tends to infinity.
\end{theorem}
\begin{proof}
At \(k=0\) the frame inequalities have no nontrivial segment.
If \(x_0=h+B_{\rm s}\), the initial presentation is already
terminal. Otherwise suppose the construction has reached \(k\)
without stopping. In particular \(x_k>h+B_{\rm s}\), the
current presentation is admitted, and its frame is weak.
Apply the first local input at this prefix. It forms the
regular coefficient from the actual companion at level \(k+1\),
using the old \(\theta_k,x_k\) and the already available
remainder data. By the second input, with
\(s_j=jc_0-\sum_{l=1}^j b_l^0\), we have \(|s_j|\le c_2\).
For each candidate segment ending at the computed value
\(x_{k+1}:=x_k-b_{k+1}\), telescoping gives
\begin{equation}\label{rt:eq-block-telescope}
 x_i-x_{k+1}-c_0(k+1-i)
     =s_i-s_{k+1}+\sum_{l=i+1}^{k+1}(f_l+q_l).
\end{equation}
The input for the last error is \(x_k\), so it obeys
\cref{rt:eq-trajectory-error-floor} before any test on
\(x_{k+1}\). The earlier errors obey the same bound by induction.
This gives the lower inequality in \cref{rt:eq-strong-drift}.
The one-step partial-sum estimate gives
\(|b_l^0|\le c_0+2c_2\). Thus
\(|b_l^0+f_l|\le C_\beta\) and \(b_l\le B_{\rm s}\);
summing gives the upper bound.

In particular \(x_{k+1}\ge x_k-B_{\rm s}>h\). The new
entry remains above the floor and extends the strong and weak
frame inequalities. Its tube deviation is zero when
\(\theta_{k+1}=x_{k+1}\). If the running minimum drops, its
drop is at most \(B_{\rm s}\). The cell condition in
\cref{rt:lem-constant-order} makes the ratio of \((\log y)^r\)
at the two successive minima less than \(L\). Taking the least
\(L\)-powers gives
\begin{equation}\label{rt:eq-adjacent-cell}
                  R_j/R_{j+1}\in\{1,L\}.
\end{equation}
The same argument with a bound \(2B_{\rm s}\) on the drop of
the minimum across two steps excludes two consecutive drops
of \(R_j\). Indeed, two drops by \(L\) would force a ratio
greater than \(L\) between the unrounded powers.
Thus the numerical frame, tube and cell requirements have been verified
for the proposed next entry. Use the one-step cofinal admission to
choose tori and cutoffs on which all reads of this finite extension are
admitted. The common-read identity preserves the coefficients already
computed when passing to a larger realization. Apply closure on each
admitted geometry to obtain the full successor presentation and its
new final remainder. The scalar update \eqref{rt:eq-vacuum-recursion}
is then evaluated for that realization. It was unnecessary for
computing the new weight. Repeating the one-step argument requires no
advance choice of the stopping horizon. On any admitted fixed finite
torus and cutoff, the same inputs realize the common coefficient
prefix through every available step.

If stopping never occurs, this procedure constructs a prefix of
every length with \(x_j>h+B_{\rm s}\). The lower drift yields
\(x_j\le\beta+c_5-\lambda_{\rm s}j\), a contradiction for
large \(j\). If \(K>0\), its preceding value is above
\(h+B_{\rm s}\), and the last decrement is at most \(B_{\rm s}\).
This gives the terminal lower endpoint. The upper endpoint
follows from the definition of \(K\). Applying the upper drift
to the whole run and the lower drift up to \(K-1\) gives
\cref{rt:eq-first-passage-bounds}. At \(K=0\), \(\beta=h+B_{\rm s}\)
and both bounds hold directly. The lower bound on \(K\)
then gives the unboundedness assertion.
\end{proof}

Only the block drift in \cref{rt:eq-strong-drift} is positive at
large separations. For example the arithmetic model
\(b_l^0=c_0+(-1)^l d\) with \(d>c_0\) has bounded one-loop
partial sums after subtracting \(c_0\) and a negative decrement
at every odd step. It satisfies the numerical premises with
vanishing corrections. Thus the stopping proof cannot require
stepwise positivity. For \(g=h^{-1/2}\) define the sets
\[
 \mathcal W_k(g)=\{\beta\ge h+B_{\rm s}:K(\beta,h)=k\}.
\]
Unbounded stopping depths make \(\mathcal W_k(g)\) nonempty
for infinitely many \(k\). The existence of one tuned value
for every prescribed depth uses a separate argument.

\rtneedspace{17\baselineskip}
\begin{proposition}[Finite realization and the torus floor]
\label{rt:prop-torus-realization}
Let the prefix of \cref{rt:thm-joint-trajectory} be given by its
common companion. If \(K'\ge K\), the numerical period conditions
at levels \(0,\ldots,K\) follow from
\[
                \ell\ge 2L^{11}M R(h+B_{\rm s}).
\]
If the stronger prospective condition stated after
\cref{rt:eq-running-cell} is required at every step, it suffices
to multiply this terminal side floor by \(L\). In a finite realization with
the supplied actual geometric and analytic admissions, the
finite construction then realizes the same weights through \(K\).
\end{proposition}
\begin{proof}
By \cref{rt:eq-adjacent-cell}, \(R_i\le L^{K-i}R_K\).
The running minimum is at most \(x_K\), so
\(R_K\le R(h+B_{\rm s})\). Since
\(P_i=\ell L^{K'-i}\ge\ell L^{K-i}\), the first
assertion follows. For the prospective condition at \(i+1\),
one needs \(P_{i+1}\ge2L^{11}M R_i\), which is obtained with
the stated extra factor. These computations select numerical
periods. To realize the known run, apply the common-read input
and the same closures on the admitted torus and cutoff. The
vacuum subtraction still uses the scalar numbers of that
realization.
\end{proof}

\section{Depth estimates and the logarithmic law}
\label{rt:sec-first-passage}

\begin{proposition}[Terminal-depth bounds]
\label{rt:prop-depth-envelopes}
In the trajectory of \cref{rt:thm-joint-trajectory}, put
\(x_{[s]}=x_{K-s}\) for \(0\le s\le K\). Then
\begin{equation}\label{rt:eq-depth-envelope}
  h+\max\{0,\lambda_{\rm s}s-c_5\}< x_{[s]}
                              \le h+B_{\rm s}(s+1).
\end{equation}
In addition,
\[
 h-c_5+\tfrac12c_0s<x_{[s]}
       \le h+B_{\rm s}+c_5+\tfrac32c_0s .
\]
At fixed \(h\), any sequence with \(K\to\infty\) has a
subsequence with a limit of \(x_{[s]}\) for every fixed \(s\).
These limits lie in the closed versions of the intervals, and
the corresponding couplings converge by continuity of
\(x\mapsto x^{-1/2}\). The limiting terminal coupling lies in
\([(h+B_{\rm s})^{-1/2},h^{-1/2}]\).

For any \(H\ge h\), the number of steps \(1\le j\le K\)
whose preceding value satisfies \(x_{j-1}\le H\) is bounded by
\(1+(H-h+c_5)/\lambda_{\rm s}\), independently of the initial
weight.
\end{proposition}
\begin{proof}
Apply \cref{rt:eq-strong-drift} from \(K-s\) to \(K\), and use
the terminal window. Each preceding value is itself above \(h\).
This gives \cref{rt:eq-depth-envelope}. For the additional
upper bound, the telescoping identity gives
\(x_{[s]}-x_K\le c_5+\tfrac54c_0s\). The lower bound follows
from \cref{rt:eq-strong-drift}. Each fixed-depth interval has
compact closure. Extract successively at depths \(0,1,\ldots\),
discarding the finitely many runs too short for the depth being
considered, and take a diagonal subsequence.

If a preceding value at most \(H\) occurs, let \(i\) be its
first index. Every counted step has preceding index between
\(i\) and \(K-1\). The lower drift from \(i\) to \(K\) gives
\(\lambda_{\rm s}(K-i)-c_5\le x_i-x_K<H-h\).
This bounds the entire remaining length, even if the weights
later rise above \(H\) again. If no such index exists, the
count is zero.
\end{proof}

\begin{lemma}[Sums of running powers]
\label{rt:lem-scale-sums}
Let \(g_j=x_j^{-1/2}\). On the prefix of
\cref{rt:thm-joint-trajectory}, for \(p>2\),
\begin{equation}\label{rt:eq-scale-power-sum}
 \sum_{j=0}^k g_j^p
 \le C_p\left(1+\frac1{\lambda_{\rm s}(p-2)}\right)
                              x_k^{1-p/2}\qquad(k\le K).
\end{equation}
For the marginal power at \(k=K\),
\[
 \frac1{B_{\rm s}}\log\!\left(1+\frac{B_{\rm s}(K+1)}
                                             {h+B_{\rm s}}\right)
 \le \sum_{j=0}^K g_j^2
 \le \frac2h+\frac1{\lambda_{\rm s}}
                 \log\!\left(1+\frac{2\lambda_{\rm s}K}{h}\right)
\]
once \(h\ge2c_5\). For every fixed \(p\ge0\) and \(\eta>0\),
let \(t_p(y)=y^{-1/2}(\log y)^p\). Above a fixed floor,
\[
  t_p(x_n)\le(1+\eta)\sqrt{n-m}\,t_p(x_m),\qquad
  t_p(x_m)\le(1+\eta)t_p(x_n)\qquad(m<n).\qedhere
\]
\end{lemma}
\begin{proof}
For \(j\le k\), the lower drift gives
\(x_j\ge x_k-c_5+\lambda_{\rm s}(k-j)\). With the fixed
floor \(x_k\ge\max\{1,2c_5\}\), this is at least
\(x_k/2+\lambda_{\rm s}(k-j)\). Bound the first term
separately and compare the remaining decreasing sequence with
the integral on the nonnegative half-line. This yields
\cref{rt:eq-scale-power-sum}. For \(p=2\) the same integral
up to \(K\) gives the upper logarithm. The upper terminal-depth
bound gives the lower logarithm by integrating the reciprocals
of \(h+B_{\rm s}(s+1)\) from zero to \(K+1\).

The profile \(t_p\) decreases above \(e^{2p}\). If \(x_m\ge
x_n\), the first ratio is at most
\((1+B_{\rm s}(n-m)/x_n)^{1/2}\), since its logarithmic factor
is at most one. It is at most \((1+\eta)\sqrt{n-m}\) above
a fixed floor. If \(x_m<x_n\), the first ratio is at most
one. In the latter case the reverse ratio is bounded by
\((1+c_5/x_m)^{1/2}\) from the lower drift, again with
logarithmic factor at most one. In the former case the reverse
ratio is at most one by monotonicity.
\end{proof}

A weighted age sum used for inherited remainder coefficients is
also a consequence. For $\kappa_0>6$ and $k\le K$,
\begin{equation}\label{rt:eq-inherited-remainder-sum}
 \sum_{j=0}^k(x_j/x_k)^2g_j^{\kappa_0}
  =x_k^{-2}\sum_{j=0}^k g_j^{\kappa_0-4}
  \le C g_k^{\kappa_0-2}.
\end{equation}
Use \eqref{rt:eq-scale-power-sum} with $p=\kappa_0-4>2$.
This estimate sums an already supplied local remainder bound;
it does not establish that bound. The positive block drift,
rather than a bound on backward excursions alone, controls the
number of contributing levels.

\begin{theorem}[Uniform ultraviolet asymptotics]
\label{rt:thm-ultraviolet-law}
Set
\[
 \omega_*(y)=\sup_{t\ge y}
                 \{\epsilon_{\rm u}(t)+\epsilon_{\rm m}(t)\},
\]
on the half-line of admitted real inverse couplings. For the
trajectory in \cref{rt:thm-joint-trajectory}, the depth error is
bounded by
\begin{equation}\label{rt:eq-depth-error}
 |x_{[s]}-x_K-c_0s|\le \Delta_h(s),\qquad
 \Delta_h(s)=c_5+\sum_{u=0}^{s-1}
       \omega_*\bigl(h+\max\{0,\lambda_{\rm s}u-c_5\}\bigr).
\end{equation}
The envelope is nondecreasing in \(s\) and
\(\Delta_h(s)/s\to0\) at fixed \(h\). In particular the
mean of \(|f_j+q_j|\) over the full run tends to zero.
At \(a_K=L^{-K}\),
for \(K\ge1\),
\[
  \left|\frac{\beta}{\log(a_K^{-1})}
                  -\frac{11}{3\pi^2}\right|
  \le \frac{h+B_{\rm s}+\Delta_h(K)}{K\log L}.
\]
In particular this convergence is uniform over the admitted
initial weights with the same stopping depth and fixed \(h\).
If the power and logarithmic envelopes of
\cref{rt:eq-trajectory-increment} hold with \(p_R>1\), then
above a fixed floor one may take
\[
     \Delta_h(s)\le C\left\{1+\sqrt{h+\lambda_{\rm s}s}\,
                         [\log(h+\lambda_{\rm s}s)]^{\mathfrak p_0}\right\}.
\]
The corresponding relative error at fixed \(h\) is
\(O_h(K^{-1/2}(\log K)^{\mathfrak p_0})\) as \(K\to\infty\).
\end{theorem}
\begin{proof}
Apply the telescoping identity on the last \(s\) steps. The
one-loop sum error is at most \(c_5\). The preceding argument
of each error in those steps has depth \(u\) between \(1\)
and \(s\). By \cref{rt:eq-depth-envelope} and the monotonicity
of \(\omega_*\), the error sum is at most the same sum at
depths \(0,\ldots,s-1\). This proves \cref{rt:eq-depth-error}.
Its summands tend to zero as the depth increases, so their
averages tend to zero. At \(s=K\), use \(x_{[K]}=\beta\)
and the terminal window, then divide by \(K\log L\).
The numerical coefficient is \(c_0/\log L\) from the
one-loop input with \(N=2\).

Above a fixed monotonicity floor the two errors have a
common majorant \(C y^{-1/2}(\log y)^{\mathfrak p_0}\).
Use \(h+\max\{0,\lambda_{\rm s}u-c_5\}\ge
(h+\lambda_{\rm s}u)/2\) for \(h\ge2c_5\).
Bound all logarithms by
\(\log(h+\lambda_{\rm s}s)\). The remaining decreasing
sum is at most
\[
 h^{-1/2}+\int_0^s (h+\lambda_{\rm s}t)^{-1/2}\,dt
 \le h^{-1/2}+2\lambda_{\rm s}^{-1}\sqrt{h+\lambda_{\rm s}s}.
\]
This proves the quantitative estimate with constants independent of
the prefix and \(h\) above the fixed floor.
\end{proof}

\begin{lemma}[Change of lattice units]\label{rt:lem-dilation}
For compatible finite lattices, let \(b\) be an integer with
\(K'-b\ge0\). In four dimensions,
\[
       Z_{K',\ell,\beta}=Z_{K'-b,L^b\ell,\beta}.
\]
At a fixed initial weight, a datum at horizon \(K'\ge K(\beta,h)\)
therefore represents the same finite Wilson integral as the
first-passage datum at side \(L^{K'-K}\ell\). The reverse change
of units is valid for the bare integral on the contracted torus.
To use the normalized history estimates there one must still
meet its geometric admission conditions.
\end{lemma}
\begin{proof}
The map \(u\mapsto L^b u\) bijects the sites of the two lattices,
their positive bonds, and their oriented plaquettes. Give
corresponding bonds the same group variables. The plaquette
products and normalized product Haar measure are unchanged,
and the Wilson action has no spacing prefactor in four
dimensions. This proves the equality. When \(K'\ge K\), take
\(b=K'-K\). Identification of intermediate renormalized
coefficients uses the common-prefix hypothesis and
\cref{rt:prop-torus-realization}, not this bare equality alone.
\end{proof}

\begin{proposition}[Common prefixes under change of units]
\label{rt:prop-common-prefix-units}
Suppose the actual supply rule is natural under the bond
relabeling of \cref{rt:lem-dilation}. The identified frame
and weight prefixes, towers, partitions and retained records
must identify the local formulas, Borel tests, measures and
determining collars. On the resulting admitted tori and
cutoffs, relabeling identifies the formed presentations
through their common prefix. For the affine reference frame
\(\theta_j=h+c_0(K'-j)\), its prefix ending at \(K'-s\) has
the same frame as the shorter run with terminal reference
value \(h_s=h+c_0s\). It is carried to the torus of side
\(L^s\ell\), for every integer \(0\le s\le K'\). A reference
side floor of the form \(\ell\ge C M R(h+B_{\rm s})\)
implies \(L^s\ell\ge C M R(h_s+B_{\rm s})\) above a fixed
floor on \(h\).
\end{proposition}
\begin{proof}
The bare measures agree by \cref{rt:lem-dilation}. At each
common step, transport the actual tests and kernels by the
same relabeling. A test indicator and its transported version
have the same value on corresponding configurations. A
draw and any projection have the corresponding pushforward
law by \cref{rt:eq-standing-kernel}. The local densities and
companions then agree on identified reads, so their coefficients
agree. Induction identifies the retained history measures
and the given local expressions at each step in the prefix.

For the affine frame,
\(h+c_0(K'-j)=h_s+c_0(K'-s-j)\) for \(j\le K'-s\).
The relabeling with \(b=s\) gives the stated torus.
For large \(y\), the ratio
\((\log(y+c_0))^r/(\log y)^r\) is less than \(L\).
Thus \(R(y+c_0)\le L R(y)\). Iterate this bound \(s\)
times at \(y=h+B_{\rm s}\) to obtain the size implication.
\end{proof}

For a self-calibrated run, a smaller terminal inverse
coupling \(h'\le h\) gives a lower stopping threshold.
If \(h'\ge Y\) and the prolonged geometries are admitted,
the same construction continues the old prefix until its
first passage at \(h'\). If the old terminal value already
lies at or below \(h'+B_{\rm s}\), the two stopping indices
agree. In either case the earlier prefix measures and
coefficients are the original ones. Comparing two different
histories or frames from a bare change of units alone
requires extra compatibility.

\section{Tuning with fixed choices}
\label{rt:sec-frozen-tuning}

Let $c_0>0$ be the one-loop drift in the coupling normalization of
this manuscript. Fix a nonnegative integer horizon $K$ and an
inverse-coupling floor $h>1$. For the slope estimates, fix an admitted
real frame satisfying, for constants $\lambda_\theta>0$ and $d_\theta\ge0$,
\begin{equation}\label{rt:eq-frozen-frame-drift}
  \theta_j\ge h,\qquad
  \theta_j-\theta_l\ge\lambda_\theta(l-j)-d_\theta
                  \quad(0\le j\le l\le K).
\end{equation}
We take $h\ge2d_\theta$ and increase its floor below. Constants
in the slope estimates may depend on the fixed drift and analytic
parameters, but not on the horizon or volume. The canonical
tuning interval will use the affine frame with $h=g^{-2}$,
\begin{equation}\label{rt:eq-frozen-frame}
       \theta_j=h+c_0(K-j),\qquad 0\le j\le K.
\end{equation}
Hold the frame and all lattice, cutoff, gauge and scaling choices
fixed as the bare weight $\beta$ varies. Write $x_j^\theta$ for
its actual weights, $x_0^\theta=\beta$, and
$b_i=x_{i-1}^\theta-x_i^\theta$ for its extracted increments. The
frame and the actual weights are different inputs to a step.

We use the coefficient comparison of \cite[Theorem A and Corollary
6.4]{EM} on the actual admitted reads. Its two estimates concern
different parameter domains. At step $i$ write
\begin{equation}\label{rt:eq-frozen-increment}
  b_i=b_i^0+f_i(z_0,\ldots,z_{i-1})+q_i(\beta),
  \qquad z_j=\theta_j-x_j^\theta(\beta).
\end{equation}
The Gaussian numbers $b_i^0$ are fixed by the reference data and
obey $|s_k|\le c_2$, where $s_k=kc_0-\sum_{i=1}^k b_i^0$.
Fix an admitted relative radius $0<\sigma_*\le1/2$. The
unmarked function $f_i$ is holomorphic on the full product
$|z_j|<\sigma_*\theta_j$ for $j<i$, with supremum at most
\[
     \epsilon_{i-1}=C_{\mathrm u}\theta_{i-1}^{-1/2}
                                  (\log\theta_{i-1})^{\mathfrak p_0}.
\]
For the bare-disc route fix $\varrho>0$, $c_\beta>0$ and $p_R>1$.
At each admitted real
centre, $q_i$ has an actual bare-variable holomorphic germ on the
disc of radius $c_\beta\varrho$ about that centre, with supremum
at most $\epsilon'_{i-1}=C_{\mathrm m}\theta_{i-1}^{-p_R}$.
This is the marked coefficient of the same actual step, including
the variation of all its supplied inherited remainder data. The
marked estimate is a separate input on its disc. All constants in
these inputs are uniform in the horizon and admissible torus. The
coefficient reads use the common companions and are real on real
data.

\begin{lemma}[A product-disc derivative bound]\label{rt:lem-polydisc-slope}
Let $f$ be holomorphic and bounded by $E$ on the finite product
$|z_j|<r_j$, with all $r_j>0$. Throughout this product,
\begin{equation}\label{rt:eq-polydisc-weighted}
 \sum_j r_j(1-|z_j/r_j|^2)|\partial_j f(z)|\le E.
\end{equation}
On the half product, $\sum_jr_j|\partial_j f(z)|\le4E/3$ and
$|f(z)-f(\widetilde z)|\le(4E/3)\max_j|z_j-\widetilde z_j|/r_j$.
The constant is independent of the number of factors.
\end{lemma}
\begin{proof}
Put $\alpha_j=z_j/r_j$. For $|t|<1$ map each coordinate by
\[
  \phi_j(t)=r_j\frac{\alpha_j+e^{i\vartheta_j}t}
                         {1+\bar\alpha_j e^{i\vartheta_j}t}.
\]
This map takes the disc into the $j$th factor and has derivative
$r_j(1-|\alpha_j|^2)e^{i\vartheta_j}$ at zero. Choose the phases
so that the nonzero derivative terms in $(f\circ\phi)'(0)$ have
the same argument. Cauchy's bound on the unit disc gives
$\sum_j r_j(1-|\alpha_j|^2)|\partial_jf(z)|\le E$. On the half
product each factor $1-|\alpha_j|^2$ is at least $3/4$. Integrating
the directional estimate on the segment in that convex product
proves the difference bound.
\end{proof}

\rtneedspace{17\baselineskip}
\begin{proposition}[Slope along a frozen prefix]
\label{rt:prop-frozen-slopes}
Use the fixed frame \eqref{rt:eq-frozen-frame-drift}. Suppose the
actual increments on a prefix satisfy \eqref{rt:eq-frozen-increment}
and the stated analytic bounds. At the real centre in question
assume $|z_j|\le\sigma_*\theta_j/2$ for all used entries. The earlier
weights and the last formed coefficient are germs of the actual
functions at this same centre. Then,
uniformly in prefix length and volume,
\begin{equation}\label{rt:eq-frozen-slope-total}
 \sum_{i=1}^k |b_i'|\le e^{A_k}(1+D_k)-1,
 \qquad
 A_k\le C h^{-1/2}(\log h)^{\mathfrak p_0},\quad
 D_k\le C\varrho^{-1} h^{1-p_R},
\end{equation}
where primes denote differentiation in $\beta$. For sufficiently
large $h$ the right side is at most $1/3$. Thus
$2/3\le (x_j^\theta)'\le4/3$ on real admitted prefixes.
On a real segment with these hypotheses throughout, the individual
increments have Lipschitz bounds $(4/3)a_i+d_i$, with $a_i,d_i$
as in \eqref{rt:eq-forced-slope}. Their sum is bounded by
$C\{h^{-1/2}(\log h)^{\mathfrak p_0}+\varrho^{-1}h^{1-p_R}\}$.
\end{proposition}
\begin{proof}
In the unmarked term, $z_j'=-(x_j^\theta)'$. For $j<i$, the frame
bound gives $\theta_j\ge\theta_{i-1}-d_\theta\ge\theta_{i-1}/2$. Lemma
\ref{rt:lem-polydisc-slope} and Cauchy's bound on the actual marked
disc give
\begin{equation}\label{rt:eq-forced-slope}
   |b_i'|\le a_i M_{i-1}+d_i,
   \qquad M_{i-1}=\max_{0\le j<i}|(x_j^\theta)'|,
   \quad a_i=\frac{8\epsilon_{i-1}}{3\sigma_*\theta_{i-1}},
   \quad d_i=\frac{2\epsilon'_{i-1}}{c_\beta\varrho}.
\end{equation}
The factor two in $d_i$ allows Cauchy circles strictly inside the
supplied disc. This differentiation includes the actual inherited
remainders through $q_i$.

Since $(x_0^\theta)'=1$, the differentiated scalar recursion bounds
$M_{i-1}$ by $V_{i-1}=1+\sum_{l<i}|b_l'|$. Then
$V_i\le(1+a_i)V_{i-1}+d_i$. Iterating gives
$V_k\le e^{A_k}(1+D_k)$ for $A_k=\sum_{i\le k}a_i$ and
$D_k=\sum_{i\le k}d_i$. For example, each product of factors
$1+a_i$ is bounded by the exponential of their sum. For $i\le k$,
\eqref{rt:eq-frozen-frame-drift} gives
\[
  \theta_{i-1}\ge h+\lambda_\theta(k-i)-d_\theta
                 \ge\tfrac12\{h+\lambda_\theta(k-i)\}.
\]
Above a fixed floor the summands decrease. Comparison with the
integrals along this affine lower envelope bounds the first sum by
$C h^{-1/2}(\log h)^{\mathfrak p_0}$ and the second by
$C\varrho^{-1}h^{1-p_R}$. For the first integral, integrate by
parts and absorb the term with exponent $\mathfrak p_0-1$ at large
$h$. The second is the integral of a power with exponent less
than $-1$. Both bounds tend to zero. Finally
$(x_j^\theta)'-1=-\sum_{i\le j}b_i'$ gives the claimed slopes.
Substitute $M_{i-1}\le4/3$ in \eqref{rt:eq-forced-slope} and
integrate over the admitted segment to obtain the Lipschitz bounds.
Their sum is at most $(4/3)A_k+D_k$.
\end{proof}

The relative disc for $f_i$ supplies one inverse power of
$\theta_{i-1}$. The bare disc for $q_i$ supplies one inverse power
of $\varrho$. The total slope error in
\eqref{rt:eq-frozen-slope-total} is
$O(h^{-1/2}(\log h)^{\mathfrak p_0}+\varrho^{-1}h^{1-p_R})$.
In particular the first term controls its order when $p_R\ge3/2$ and
$\varrho$ is fixed. The derivative bound for a product disc gives no
decay with the age of a particular input coordinate.

\begin{proposition}[Slope with an independent marked polydisc]
\label{rt:prop-independent-marked-slope}
Use the fixed frame \eqref{rt:eq-frozen-frame-drift} and the
stated unmarked coefficient bound. Suppose instead that the actual
marked coefficient factors through the earlier deviations in
\eqref{rt:eq-frozen-increment} as a
holomorphic function $q_i(z_0,\ldots,z_{i-1})$. Assume that it has
supremum at most $C\theta_i^{-p_R}$ on the product of discs of radii
$c_{\rm m}\theta_j(\log\theta_j)^{-C_*}$, $j<i$, for fixed $c_{\rm m}>0$,
$C_*\ge0$ and $p_R>0$. This factorization includes the genuine
inherited dependence. Assume the runs remain in the half products
for both terms, with the same compatible core and coefficient
choices. Then the bare slope satisfies
\begin{equation}\label{rt:eq-independent-marked-slope}
   \sum_{i\le k}|b_i'|\le e^{A_k^\flat}-1,
   \qquad
   A_k^\flat\le C\bigl[h^{-1/2}(\log h)^{\mathfrak p_0}
                         +h^{-p_R}(\log h)^{C_*}\bigr].
\end{equation}
The bound is uniform in $K$, $k$ and compatible volume.
\end{proposition}
\begin{proof}
Above a fixed floor $y(\log y)^{-C_*}$ increases. For $j<i$,
the frame bound gives $\theta_j\ge\theta_i-d_\theta\ge\theta_i/2$.
The marked radii are therefore bounded below by a fixed multiple
of $\theta_i(\log\theta_i)^{-C_*}$. The unmarked radii are bounded
below by $\sigma_*\theta_{i-1}/2$, and $\theta_{i-1}\ge\theta_i/2$.
Using eventual monotonicity of $y^{-3/2}(\log y)^{\mathfrak p_0}$
in the unmarked bound, \cref{rt:lem-polydisc-slope} gives the
absolute derivative row bound
\[
  a_i^\flat\le C\bigl[\theta_i^{-3/2}(\log\theta_i)^{\mathfrak p_0}
                   +\theta_i^{-p_R-1}(\log\theta_i)^{C_*}\bigr].
\]
The chain rule gives $|b_i'|\le a_i^\flat M_{i-1}$. The scalar
iteration in the preceding proof with $d_i=0$ gives
$1+\sum_{i\le k}|b_i'|\le\exp(\sum_{i\le k}a_i^\flat)$. Here
$\theta_i\ge\{h+\lambda_\theta(k-i)\}/2$. Compare both sums with
the integrals of their decreasing envelopes above $h/2$.
Writing $y=(h/2)u$ leaves an integrable power $u^{-1-\delta}$
times a fixed power of $1+\log u$, with $\delta>0$. This gives
\eqref{rt:eq-independent-marked-slope}.
\end{proof}

\subsection{The interval of admitted bare weights}

Let $\mathsf t^\sharp(y)=\varpi y(\log y)^{-2q_{\mathrm t}}$ with
$\varpi>0$ and $q_{\mathrm t}\ge0$. A tube is a positive
nonincreasing sequence $\mathsf t_j$, with
$\mathsf t_j\le\mathsf t^\sharp(\theta_{j+1})$ for $j<K$.
For $K=0$ use $\mathsf t_0\le2\mathsf t^\sharp(\theta_0)$.
For any frame of length $K\ge1$ the largest tube has entry
$\min_{1\le i\le\min(j+1,K)}\mathsf t^\sharp(\theta_i)$ at $j$.
The preceding upper bounds and monotonicity require this minimum,
and these minima themselves satisfy the requirements. A constant
tube of radius $t$ is permitted whenever
$t\le\varpi x_*^{3/4}/C_a$ and every frame entry is at least $x_*>e$,
where $a=2q_{\mathrm t}$ and $C_a=\max\{1,(4a/e)^a\}$ with
$C_0=1$ in this formula. Indeed, maximizing
$u^a e^{-u/4}$ gives $(\log y)^a\le C_a y^{1/4}$ for $y\ge1$.

For the canonical interval construction, now specialize to the
affine frame \eqref{rt:eq-frozen-frame}. Put $\iota_j=\epsilon_j+\epsilon'_j$ and choose
\begin{align}
 T_K&=\varrho+c_2,&
 T_j&=T_K+\sum_{l=j}^{K-1}\iota_l,&
 \widehat{\mathsf t}_j&=T_j+c_2+8\varrho,\label{rt:eq-canonical-tube}\\
 F_j(\beta)&=\theta_j-x_j^\theta(\beta)+s_j.\nonumber
\end{align}
Here $F_j$ is defined once the weight at $j$ is computed, before
requiring closure at that weight. The sum of $\epsilon_l$ in
$T_j$ is at most $C\theta_j^{1/2}(\log\theta_j)^{\mathfrak p_0}$.
To see this, bound its logarithm factor by the one at $\theta_j$
and sum the inverse square roots on the reference line by an
integral. The sum of $\epsilon'_l$ is at most $C h^{1-p_R}$.
For $j<K$ and sufficiently large $h$, $\theta_{j+1}\ge\theta_j/2$.
Consequently the ratio of $\widehat{\mathsf t}_j$ to
$\mathsf t^\sharp(\theta_{j+1})$ tends uniformly to zero as $h$
increases. The terminal value of $\widehat{\mathsf t}$ is fixed,
and the one-entry convention holds at the same large floor.
The sequence in \eqref{rt:eq-canonical-tube} is therefore a tube.
Increasing the floor also makes
$\widehat{\mathsf t}_j\le\sigma_*\theta_j/2$ at every level.
At retained entries $x_j^\theta\ge\theta_j-T_j-c_2\ge\theta_j/2$.
Choose $h/2$ above all absolute input floors. The terminal range
$[h-\varrho,h+\varrho]$ then also stays above those floors.

\rtneedspace{27\baselineskip}
\begin{theorem}[Full tuning in the canonical interval]
\label{rt:thm-frozen-tuning}
Use the reference frame \eqref{rt:eq-frozen-frame}, the actual
increment bounds of \eqref{rt:eq-frozen-increment}, and the uniform
one-loop partial-sum bound above. Assume the following admitted
prefix operation. The initial Wilson inputs are admitted on the
interval $|\theta_0-\beta|\le T_0$. From each real prefix with
$|F_j|\le T_j$ at every closed level, the actual next coefficient
is formed on a neighbourhood before imposing the new bound on
$F_{j+1}$. The coefficient and all earlier weights there have the
analytic inputs of \cref{rt:prop-frozen-slopes}. When the computed
entry satisfies $|F_{j+1}|\le T_{j+1}$, the actual next level closes
with the declared tube room. This operation includes the inherited
history and vacuum data at each closed prefix. It uses one frame,
with common companions and coefficients on the admissible tori.

For a floor on $h$ independent of $K$, there are nested compact
intervals $I_j$ of actual closed prefixes with
$|F_l|\le T_l$ for $l\le j$. On $I_j$, $F_j$ decreases from $T_j$
to $-T_j$. For each real $v$ with $|v-h|\le\varrho$ there is one
$\beta\in I_K$ such that $x_K^\theta(\beta)=v$. The tuned bare
weight $\beta_K^\theta$ for $v=h$ satisfies
\begin{equation}\label{rt:eq-full-tuned-weight}
 \tfrac12c_0 K-2c_2\le \beta_K^\theta-h
                   \le\tfrac32c_0 K+2c_2,
 \qquad
 |\beta_K^\theta-h-c_0K|\le c_2+\sum_{j<K}\iota_j .
\end{equation}
For all $\beta\in I_K$ and $0\le i\le j\le K$ one has
\[
 \tfrac12c_0(j-i)-2c_2\le x_i^\theta-x_j^\theta
                   \le \tfrac32c_0(j-i)+2c_2 .
\]
The uniqueness assertion concerns the fixed frame and its
canonical interval.
\end{theorem}
\begin{proof}
The initial $F_0=\theta_0-\beta$ maps $I_0$ decreasingly onto
$[-T_0,T_0]$. The high floor makes all initial weights positive.
At a closed prefix the bound $|F_l|\le T_l$ gives
$|z_l|\le T_l+c_2$, leaving room $8\varrho$ inside the chosen tube.
The premise supplies the next coefficient before any closure at
the next weight. Its real function $F_j$ on a neighbourhood of
$I_{j-1}$ has derivative $-(x_j^\theta)'\le-2/3$ by
\cref{rt:prop-frozen-slopes}. Moreover
\begin{equation}\label{rt:eq-frozen-F-difference}
     F_j-F_{j-1}=b_j-b_j^0,
     \qquad |F_j-F_{j-1}|\le\iota_{j-1}.
\end{equation}
This follows directly from the changes in the frame and in $s_j$.
At the left endpoint of $I_{j-1}$, $F_{j-1}=T_j+\iota_{j-1}$,
so $F_j\ge T_j$. At the right endpoint $F_j\le-T_j$. The strict
slope and continuity give the unique ordered preimages of $T_j$
and $-T_j$. Their interval $I_j$ meets the new bound at every point.
The supplied closure operation applies on it with the same room.

At $K$, $x_K^\theta=h+s_K-F_K$. Its image contains
$[h-\varrho,h+\varrho]$ because $T_K=\varrho+c_2$ and
$|s_K|\le c_2$. Strict monotonicity gives the stated uniqueness.
Summing the increments from $i+1$ through $j$ gives
\[
  x_i^\theta-x_j^\theta=c_0(j-i)+s_i-s_j
                                  +\sum_{i<l\le j}(b_l-b_l^0).
\]
The high floor gives $\sup_l\iota_l\le c_0/2$. This proves the
spacing bounds. At the tuned value $x_K^\theta=h$ and $s_0=0$,
the same sum gives \eqref{rt:eq-full-tuned-weight}. For $K=0$,
the initial interval and $x_0^\theta=\beta$ already prove all
assertions.
\end{proof}

On $I_K$ the terminal slope gives, for a bare offset $t$ with
$\beta_K^\theta-t\in I_K$,
\[
       x_K^\theta(\beta_K^\theta-t)=h-t+O(\eta_h|t|),
       \qquad \eta_h=e^{A_K}(1+D_K)-1.
\]
This follows by integrating the slope error on the intervening
real segment. The tuned endpoint $h$ is a condition on the fixed
reference run. First passage is defined separately for the
self-calibrated trajectory.

\section{A backward equation for homogeneous profiles}
\label{rt:sec-homogeneous-tuning}

Use the complete profile spaces and the scalar interface of
\cite[Theorem B and Corollary~7.5]{EM}. For a prescribed weight
sequence $x$, the forward profile $\psi_j(x)$ starts at zero.
At each step the actual homogeneous formation carries the scalar
$x_j-b^{\mathrm G}_{j+1}-r_j(x_j,\psi_j(x))$. Here
$b^{\mathrm G}_{j+1}$ is the companion coefficient of the actual
reference Gaussian, and $r_j$ is the bounded companion readout of
the actual unprojected formation. The Gaussian coefficient is fixed
when the action weights vary. For an endpoint reference $X$ at level
$n$, the nominal frame at level $j$ is
$a_j=X+\sum_{\ell=j}^{n-1}b^{\mathrm G}_{\ell+1}$. On the declared
relative weight box about this frame the forward result supplies
\begin{equation}\label{rt:eq-homogeneous-interface}
 \begin{aligned}
 |r_j(x_j,\psi_j(x))|&\le C e(a_j),\\
 |r_j(x_j,\psi_j(x))-r_j(\widetilde x_j,\psi_j(\widetilde x))|
     &\le C k(a_j)\|x-\widetilde x\|_\infty,\\
 \|\psi_j(x)-\psi_j(\widetilde x)\|_{j,*}
     &\le C k(a_j)\|x-\widetilde x\|_\infty.
 \end{aligned}
\end{equation}
The functions are $e(y)=y^{-1}(\log y)^m$ and $k(y)=e(y)/y$
with fixed $m>1$. The first two bounds hold for $0\le j<n$ and
concern the unprojected scalar interface. The profile difference
bound holds for $0\le j\le n$, including the terminal profile.
These are inputs to this construction.

\rtneedspace{25\baselineskip}
\begin{theorem}[Homogeneous terminal tuning]\label{rt:thm-homogeneous-tuning}
Fix a horizon $n\ge0$, a positive real endpoint reference $X$, and
real Gaussian coefficients with
\[
   \left|\sum_{\ell=j}^{n-1}(b^{\mathrm G}_{\ell+1}-b_\infty)\right|
             \le B_a,\qquad 0\le j\le n,
\]
for fixed $b_\infty>0$ and $B_a<\infty$. Define
\[
     a_j=X+\sum_{\ell=j}^{n-1}b^{\mathrm G}_{\ell+1},
     \qquad a_n=X.
\]
Assume the admitted actual forward profile and scalar operations of
\eqref{rt:eq-homogeneous-interface} on the frozen frame $a$, with
holomorphy and the uniform profile bound on the buffered domains of
the forward theorem. Their exact physical exponent includes the
scalar vacuum factor as a separate number. Fix a relative radius
$\delta\in(0,1/2]$ within the inner profile weight domain.

Above a floor on $X$ independent of the horizon, for every complex
$\zeta$ with $|\zeta|<cX$ for a fixed $c>0$, there is one trajectory
in the box $|x_j-a_j|\le\delta a_j$ with initial profile zero and
terminal weight $x_n=X+\zeta$. It is real for real data and real
$\zeta$. The trajectory and its forward profiles are holomorphic in
$\zeta$. At $\zeta=0$, $x_j\asymp X+n-j$ and
\begin{equation}\label{rt:eq-homogeneous-tuned-law}
     x_0=X+b_\infty n+O\bigl(1+(\log(X+n))^{m+1}\bigr).
\end{equation}
All these constants are uniform in $n$ and the admitted volume. On
a fixed smaller endpoint disc,
\begin{equation}\label{rt:eq-homogeneous-endpoint-jets}
 \sup_j|x_j'-1|\le C\frac{(\log X)^m}{X},\qquad
 \|\psi_j'\|_{j,*}\le C k(a_j),\qquad
 |x_j^{(s)}|\le C_s\frac{(\log X)^m}{X^s}\quad(s\ge2).
\end{equation}
In particular the real bare weight is strictly increasing with the
real endpoint there.
\end{theorem}
\begin{proof}
The Gaussian bound gives $|a_j-X-b_\infty(n-j)|\le B_a$. In
particular $a_j\asymp X+n-j$ above a fixed floor, even if some
Gaussian increments are negative. Compare the sums with those
over $X+c_1 q$ for $q\ge0$, using fixed $c_1>0$. The first
comparison integrates $(\log y)^m/y$ up to $C a_j$. The second
integrates $(\log y)^m/y^2$ on a half-line, where integration by
parts bounds it by $C(\log X)^m/X$. Including the first terms gives
\begin{equation}\label{rt:eq-homogeneous-sums}
 \sum_{\ell=j}^{n-1}e(a_\ell)\le C\bigl(1+(\log a_j)^{m+1}\bigr),
 \qquad
 \sum_{\ell<n}k(a_\ell)\le C(\log X)^m/X.
\end{equation}

The backward sum of the exact outgoing weights is the fixed-point
equation
\begin{equation}\label{rt:eq-homogeneous-backward}
   (\Phi^\zeta x)_j=a_j+\zeta+
             \sum_{\ell=j}^{n-1}r_\ell(x_\ell,\psi_\ell(x)).
\end{equation}
Take $c\le\delta/8$. Since $a_j\ge X/2$ for large $X$, the endpoint
term is at most $\delta a_j/4$. By the first bound in
\eqref{rt:eq-homogeneous-sums}, increasing the floor makes the
remainder sum at most $\delta a_j/4$. Thus the image is within the
half box. The second bound in \eqref{rt:eq-homogeneous-sums} and
the scalar difference bound in \eqref{rt:eq-homogeneous-interface}
give a contraction in the unweighted supremum norm with rate
$\mathfrak k_X\le C(\log X)^m/X$. Increase the floor so that
$\mathfrak k_X\le1/5$. The closed finite box is complete. Iteration
gives one fixed point, and taking adjacent differences in
\eqref{rt:eq-homogeneous-backward} gives the actual outgoing-weight
recursion. Every trajectory in the box has the prescribed forward
profiles and must satisfy that equation, which proves uniqueness
there.

Iteration starting at $a$ preserves reality on the real interval.
For complex $\zeta$ the iterates are holomorphic because the
prescribed profiles and scalar readout are holomorphic on the
larger admitted domain. The uniform contraction gives normal
convergence on smaller endpoint discs, so the fixed point is
holomorphic there. Summing the scalar correction and using
\eqref{rt:eq-homogeneous-sums} gives \eqref{rt:eq-homogeneous-tuned-law}.
The exact unprojected step identifies the resulting physical
profile and carried weight at each level with its actual integral.
The separate scalar vacuum factor is the one of that step.

At the fixed point, which is in the interior of the box, the
operator norm of $D_x\Phi^\zeta$ is at most $\mathfrak k_X$.
Differentiating the fixed-point equation gives
$\|x'\|_\infty\le1+\mathfrak k_X\|x'\|_\infty$ and hence
$\|x'-1\|_\infty\le\mathfrak k_X/(1-\mathfrak k_X)$.
The forward profile estimate gives the claimed bound on $\psi_j'$.
Apply Cauchy's estimate to $x_j'-1$ on circles of radius a fixed
multiple of $X$ inside the larger endpoint disc. Each additional
derivative costs a factor $X^{-1}$, giving the higher jets in
\eqref{rt:eq-homogeneous-endpoint-jets}. The real slopes lie in
$[3/4,5/4]$. When $n=0$, the sums are empty and the equation is
$x_0=X+\zeta$, with zero profile.
\end{proof}

The actual operations in this theorem are the homogeneous regular
ones on their fixed physical radii. Their scalar interface is
separate from the actual marked bare-disc dependence used in
\cref{rt:thm-frozen-tuning}.

\section{Comparison of admitted presentations}
\label{rt:sec-presentation-matching}

For a real run over a fixed frame, let $x_j$ be the weight used
in the Gaussian presentation. A uniform scalar shift $\mathsf c_j$ is
carried in the local action as $\mathsf c_j A$. The effective Wilson
weight in the all-small-field exponent is $x_j-\mathsf c_j$. We consider
presentations with one Euclidean regular member at each level,
one companion coefficient and one scalar shift, with the real
weights and shifts admitted in their stated domains. All tests,
charts and underlying integration choices are those of the common
frame. In particular the shift is not a spatially varying source.

Write $b_j=x_{j-1}-x_j$ for the chosen decrement, and $\widehat b_j$
for the scalar actually read from the new regular member. The
bookkeeping identities are
\begin{equation}\label{rt:eq-effective-weight}
 \begin{aligned}
  \mathsf c_j&=\mathsf c_0+\sum_{l=1}^j(\widehat b_l-b_l),\\
  x_j-\mathsf c_j&=x_{j-1}-\mathsf c_{j-1}-\widehat b_j .
 \end{aligned}
\end{equation}
The first line follows by successively equating the
effective outgoing coefficient to the incoming coefficient minus
the coefficient read from the physical step.

\begin{proposition}[Matching on common all-small-field windows]
\label{rt:prop-presentation-matching}
Consider two such real presentations, with primed symbols for the
second, on one fixed frame. Assume their unnormalized all-small-field
measures agree on every admitted torus through level $n$. At each
level $m$ assume an open physical window $\mathfrak O_m$ about the
identity where the relevant tests equal one and the physical density
has its positive continuous all-small-field formula. These windows
are hereditary. The earlier averages of a background in
$\mathfrak O_m$ belong to the corresponding windows. The complete
retained physical pairs, including their physical currents, agree
at the corresponding earlier reads up to the admitted site gauges. For each outer datum, use the common parametrization of the actual
free fluctuations on a connected complex domain containing the real
cut and zero fluctuation. The incoming physical exponent and each omitted remainder total,
composed with their actual physical reads, are holomorphic there.
A nonempty open set in the full real fluctuation coordinate space
maps all those reads into the hereditary windows. These are
hypotheses on the actual fluctuation maps, not on a detached
independent-current extension.

Assume the exact common formation and vacuum conventions. The
centered raw regular integral uses the incoming effective exponent,
with the omitted earlier remainder terms removed by their actual
common reads. The two Gaussian presentations of its logarithm
separate the same reference constants and the respective term
$-(D_j/2)\log(x_{j-1}/\theta_{j-1})$ at step $j$, where $D_j$ is
the real dimension of the free fluctuation coordinates. Write
$\mathscr E_l(P)=\sum_{X,z}\mathcal E_l(X,P,z)$ and
$\mathscr R_l(P)=\sum_X\mathcal R_l(X,P)$ for the complete raw totals
on the physical pair $P$. Let $P_l(W)$ be the retained physical pair
at the prescribed earlier averaged datum, on which the level-$l$
family is read in the window $\mathfrak O_m$.
On $\mathfrak O_m$ the exact
raw all-small-field density is taken in a common reference measure
with full support on each open window, such as product Haar
measure in the free link coordinates. It has the form
\begin{equation}\label{rt:eq-matching-density}
  p_m(W)\exp\left\{E_{(m)}-(x_m-\mathsf c_m) A(U_m(W))
      +\sum_{l\le m}\bigl(\overline{\mathscr E}_l(P_l(W))
                       +\overline{\mathscr R}_l(P_l(W))\bigr)\right\}.
\end{equation}
Here $p_m$ is the same positive prefactor in both presentations,
$\overline{\mathscr E}_l$ subtracts the trivial value and its own
extracted classical term, and
$\overline{\mathscr R}_l=\mathscr R_l-\mathscr R_l(P_\circ)$.
The trivial datum is $P_\circ=(\mathbf1,0)$. Require
$P_l(\mathbf1)=P_\circ$ up to the admitted gauge for every retained
read. In particular the actual physical current vanishes at the
trivial field, and the trivial background satisfies $A(U_m(\mathbf1))=0$.
The complete remainder at $m$ is used after closure at $m$.
The once-counted increment of $E_{(m)}$ in
\eqref{rt:eq-vacuum-recursion} contains the displayed Gaussian
scalar and $\mathscr E_m(P_\circ)+\mathscr R_m(P_\circ)$.
Its other factors are the same in both presentations. Use cofinal
compatible tori and the common physical coefficient readout of
\cite[Proposition~6.3]{EM}.

Then $x_j-\mathsf c_j=x'_j-\mathsf c'_j$ through $n$. The extracted coefficients
match, and on every corresponding hereditary background the
regular totals obey
\begin{equation}\label{rt:eq-regular-presentation-constant}
  \mathscr E_j(P)-\mathscr E'_j(P)
            =\frac{D_j}{2}\log(x_{j-1}/x'_{j-1}).
\end{equation}
The remainder totals and the accumulated vacuum agree through $n$
on their common all-small-field windows. If the next regular formation
satisfies the same common-read, cut and continuation hypotheses on
a common physical window, its coefficient agrees as well.
\end{proposition}
\begin{proof}
At the bare level equality of the positive densities identifies the
effective Wilson weight. The action is zero at the identity and
is nonconstant on its neighbourhood. On a later open window equality
of the measures gives pointwise equality of the continuous densities.
Indeed the stipulated common reference measure has full support there. Taking the real
logarithm introduces no branch ambiguity.

Proceed by induction. The preceding regular and remainder identities
hold on the earlier backgrounds of the new window by heredity. If
the formation omits recent remainders from its regular list, their
totals already agree at the actual common reads. The composed differences vanish on the stated real open fluctuation
patch. Successive one-variable identity theorems make them zero
on a complex neighborhood, and connectedness makes them zero
throughout the common fluctuation domain. Thus the incoming
physical exponents in the next centered regular integral agree.
The possibly Borel cut indicators are applied only on the real slice. Its subtraction at zero fluctuation
is also identical. Integration on the same real cut gives the same
raw regular logarithm. The prescribed scalar difference gives
\eqref{rt:eq-regular-presentation-constant}. The curvature coefficient
of a field-independent total is zero. The cofinal companion readout
therefore gives $\widehat b_j=\widehat b'_j$, and
\eqref{rt:eq-effective-weight} identifies the new effective weights.

The regular difference in \eqref{rt:eq-regular-presentation-constant}
cancels after its value at the trivial datum is subtracted. The
classical terms agree by equality of the extracted coefficients.
At $W=\mathbf1$ the right side of \eqref{rt:eq-matching-density}
is $p_m(\mathbf1)e^{E_{(m)}}$. Equality there identifies the accumulated
vacuum. Dividing by the common effective Wilson, regular and
earlier remainder factors identifies the vacuum-subtracted new
remainder total. In the prescribed vacuum increment, the change of the
Gaussian scalar cancels the change of the regular trivial value
by \eqref{rt:eq-regular-presentation-constant}, and the other factors
are common. Equality of the vacuum increments therefore gives
equality of the new remainder trivial values. This finishes the
induction. At a next regular formation the comparison through the
raw regular integral and coefficient uses only the preceding equal
data. It does not require closure at the new weight.
\end{proof}

The use of this result for a constant scalar shift is precise. Take
an admitted bare centre $\beta$ and a real shift $\zeta$ on a common
admitted interval. Compare the presentation starting at
$\beta-\zeta$ with zero initial shift and the presentation using
the Gaussian weights $x_j^\theta(\beta)$ of the admitted run at
$\beta$, with initial shift $\zeta$ carried subsequently as
$\zeta_j$. If their actual all-small-field measures and reads meet
the hypotheses above, then
\begin{equation}\label{rt:eq-source-shift-match}
          x_j^\theta(\beta-\zeta)=x_j^\theta(\beta)-\zeta_j.
\end{equation}
The two initial densities coincide since
$e^{-\beta A+\zeta A}=e^{-(\beta-\zeta)A}$. When both presentations
use the same fixed labelled kernels, each all-small-field measure
is obtained by applying the same linear kernel to this density.
This supplies the measure equality for the application. It uses
the actual identity of the tests, gauge weights and retained maps
in each labelled kernel. On an admitted real segment with the
stated slope hypotheses, \cref{rt:prop-frozen-slopes} gives
$|\zeta_j|\le4|\zeta|/3$. The fixed frame may be the affine
reference frame or the weight prefix of the self-calibrated run
at the bare centre. In the latter case hold that whole reference
prefix fixed while varying the bare weight on the segment.
Differentiating \eqref{rt:eq-source-shift-match} at $\zeta=0$
gives $\partial_\zeta\zeta_j|_{\zeta=0}=(x_j^\theta)'(\beta)$.
This is the normalization of the running constant insertion in this
comparison. Identification with any separate observable construction
requires that construction's own actual matching.

The bare-variable holomorphic germs about centres of a canonical
interval $I_k$ can be glued when they agree with the actual real
run on the real part of that interval. For two distinct centres in
$I_k$, intersection of their open discs implies that their common
real intervals contain an open segment between the centres. For
coincident centres, use the one-sided interval inside $I_k$ if
necessary. The identity theorem on the convex disc intersection
gives equality there. For regular and remainder totals, first fix
a physical datum in a window common to both germs throughout an
open admitted real subinterval of the overlap. Their actual totals
agree there, with the fixed scalar convention, so the same argument
identifies their bare-variable continuations at that datum on the
connected common domain containing that real subinterval. This
identification uses the stated common physical window. A comparison
at centres from different real components
needs a common admitted real interval or a separately supplied
identification of the analytic branches. This gluing uses one frame
and compares complete physical totals on the stated windows.

\rtneedspace{19\baselineskip}
\begin{lemma}[A weighted two-run estimate]\label{rt:lem-weighted-two-run}
Let $r_0,n$ be integers with $0\le r_0<n$, let $0<\rho<1$ and $\Pi\ge0$, and set
\[
    S_w=\frac{\rho^{n-w+1}}{1-\rho}+(w-r_0)\Pi,
    \qquad r_0<w\le n.
\]
Suppose the nonnegative quantities $y_w,D_u,k_{u,w},h_u,h'$ of
a finite comparison satisfy, for $r_0\le s,u<n$,
\begin{align*}
 y_{s+1}&\le\sum_{u=r_0}^sD_u, &
 D_u&\le M_*(\rho^{n-u}+\Pi)+\sum_{w=u+1}^n k_{u,w}y_w,\\
 \sum_{w=u+1}^n k_{u,w}S_w&\le h_u S_{u+1}+h'\Pi, &
 \sum_{u=r_0}^{n-1}h_u+h'&\le1/2,
\end{align*}
where $M_*\ge0$. Then $y_w\le2M_*S_w$ for every indicated $w$.
\end{lemma}
\begin{proof}
The weights $S_w$ are positive and nondecreasing. Put
$Z=\max_{r_0<w\le n}y_w/S_w$. The forcing sum through $s$ is at
most $M_*S_{s+1}$ by the geometric sum. The other sum is bounded by
\[
 Z\left(\sum_{u=r_0}^s h_u S_{u+1}+(s-r_0+1)h'\Pi\right)
          \le\tfrac12 ZS_{s+1}.
\]
Divide by $S_{s+1}$ and take the maximum. This gives $Z\le2M_*$.
\end{proof}

The kernel inequality in this last lemma is a premise on the
actual comparison of two runs. It allows dependence on later
depths. A bound on the single-run increments alone supplies no
such kernel estimate.

\section{The zero-source partition function}
\label{rt:sec-partition}

Fix an admitted real run through level \(n\). Its terminal measure is
\(\mu_n\), and the common subtraction \(E_{(n)}\) is given by
\eqref{rt:eq-vacuum-recursion}. Let \(\mathcal S\) be a standing
label at level \(n\). Histories with that label have already been
summed in its measure \(\tau_{\mathcal S}\). In particular the
empty label includes histories whose large-field components have
been healed. Positivity and exactness give
\begin{equation}\label{rt:eq-history-decomposition}
   \mu_n=\sum_{\mathcal S}\tau_{\mathcal S},
   \qquad \tau_{\mathcal S}\ge 0,
   \qquad \sum_{\mathcal S}\tau_{\mathcal S}(1)=Z_{K',\ell,\beta}.
\end{equation}
Each measure can be projected to the final links for the purposes
of this section. The history sum does not require an exponential
formula for the measure of a single raw history.

\subsection{Bounding the retained integrals}
\label{rt:subsec-retained-integrals}

The exponent of a label has a numerical part \(v_{n,\mathcal S}\).
Use the sign in which the numerical factor in the measure is
\(\exp v_{n,\mathcal S}\). For the empty label this factor is
\(\exp E_{(n)}\). For another label the difference
\(v_{n,\mathcal S}-E_{(n)}\) stays in its normalized contribution.
The absolute value of this difference need not be bounded by a
constant per final site if an old fine region is present.

\begin{lemma}[Majorizing coupled integrals]\label{rt:lem-component-majorant}
Let \(\xi_1,\ldots,\xi_b\) be disjoint groups of retained
variables, with sigma-finite reference measures and product
reference measure \(\prod_i\mu_i(d\xi_i)\). Let \(y\) denote the remaining variables,
with a probability reference measure \(\pi(dy)\). Suppose the
factor \(f_i(\xi_i,e_i)\) is evaluated at an admissible exterior
argument \(e_i=e_i(y,\xi_1,\ldots,\xi_b)\). Assume the admitted
support is Borel and the evaluated factors are measurable there.
Include the support restrictions in the domain of integration.
If \(\mathcal D_i(\xi_i)\)
contains all those exterior arguments, then an integral of the
absolute product of the factors is bounded by
\begin{equation}\label{rt:eq-retained-product}
   \int \prod_i |f_i(\xi_i,e_i)|\,d\pi(y)\prod_i d\mu_i(\xi_i)
   \ \le\
   \prod_i \int^*\sup_{e\in\mathcal D_i(\xi_i)}
                         |f_i(\xi_i,e)|\,d\mu_i(\xi_i).
\end{equation}
Assume each upper integral on the right is finite. The integral
on the left is over the admitted support. Here
\(\int^*\) is the upper integral over nonnegative measurable
majorants. The same conclusion holds with a nonnegative Borel factor
depending only on \(\xi_i\) inserted on both sides for each \(i\),
provided the modified upper integrals are finite.
\end{lemma}
\begin{proof}
For each \(i\) choose a measurable majorant of its supremum, taking
the supremum over an empty exterior set to be zero. Drop the support
indicator after applying the majorants. The product then depends
only on the disjoint groups \(\xi_i\), so Tonelli's theorem gives the
product of their integrals. Integrating \(y\) costs one. Taking
majorants with integrals arbitrarily close to the finite upper
integrals proves the estimate used here. The inserted factors can
be included in the definition of \(f_i\).
\end{proof}

The same estimate controls composition of retained kernels.
For a measurable kernel $K(\xi;p)$, nonnegative Borel site weight
$v(\xi)$ and $c\ge0$, set
\[
 \|K\|_c=\int^*e^{c v(\xi)}\sup_p|K(\xi;p)|\,d\mu(\xi).
\]
If this is finite, the integral operator
$(T_KF)(p)=\int K(\xi;p)F(\xi,p)\,d\mu(\xi)$ satisfies
$\|T_K F\|_\infty\le\|K\|_0\|F\|_\infty$ for bounded Borel $F$.
For two kernels, suppose the actual measurable read $q(\xi,p)$
belongs to the second exterior domain at $\eta$ for every
$(\xi,\eta,p)$ on the admitted joint support. Extend the evaluated
integrand by zero outside that support.
Bound the joint supremum pointwise by the two local suprema,
then use the product-majorant argument of
\eqref{rt:eq-retained-product} to obtain
\[
 \int^* e^{c(v_1(\xi)+v_2(\eta))}
  \sup_p |K_1(\xi;p)K_2(\eta;q(\xi,p))|
                 \,d\mu_1(\xi)d\mu_2(\eta)
       \le \|K_1\|_c\|K_2\|_c.
\]
Both norms are assumed finite. Discrete record sums with their
normalized frame weights can be included in the reference measures.
The actual read inclusion is required before making this estimate.

The distinction between the two orders of supremum and integration
is necessary. For \(0<\epsilon<1\), on \([0,1]^2\) take both
component factors to be \(1_{\{|u-v|<\epsilon\}}\). At each fixed
exterior value the integral in the component's own variable is at
most \(2\epsilon\). Their product has integral
\(2\epsilon-\epsilon^2\), which exceeds \(4\epsilon^2\) for
\(\epsilon<2/5\). The upper integrals in
\eqref{rt:eq-retained-product} do bound this product.

For a retained operation born on \(X\) at level \(i\), write
\(\omega\) for its discrete record, \(\xi\) for its stored
variables and \(\mu_\omega\) for their reference measure. On its
support let \(w_\omega(\xi,e)\) be the actual integrand with the
energy factors assigned to this operation. Let \(\varpi(\omega)\)
be the product of the normalized reference-frame weights for this
record. Other discrete choices have counting weight one. The
supremum in its weighted mass is taken in a determining domain
chosen at its birth using only the prefix through level $i$.
The supplied invariant inclusion places every later admitted
evaluation in that birth domain. It is not enlarged after a
future evaluation has been selected.
In a drop-out it is taken over the admitted real evaluations
compatible with the stored variables and the recorded tests,
except for the test that caused the drop-out. The mass used in the
comparison is
\begin{equation}\label{rt:eq-weighted-retained-mass}
   \mathsf M_c(i,X;\mathcal B)=
   \sum_{\omega\in\mathcal B} \varpi(\omega)
      \int^* e^{c\,v(\omega)}
         \sup_e |w_\omega(\xi,e)|\,d\mu_\omega(\xi).
\end{equation}
Here \(\mathcal B\) can restrict the visible records of the
operation, and \(v(\omega)\) is the sum of
\(1+\log x_t\) over its charged stored sites at their own levels
\(t\). In this sum each site at a given level is charged at most a
fixed number of times. A mass bound is required for the actual
reference measures, including their Jacobians or densities when
these occur in the step.

An island includes its assigned core Wilson factor. A drop-out
uses a share of the Wilson action of the term that carries it,
over bare plaquettes in its assigned core, because the retained
operation by itself has no such factor. A failure of a plaquette
test uses an assigned positive Wilson contribution; a failure of
an approximate fluctuation test may also use a retained gauge
contribution. The genuine fluctuation, collar and rim tests use
the corresponding positive quadratic contributions on their
carriers. More precisely, if a certified failure forces an
allocated nonnegative energy \(Q\ge 2a\) on the test support,
then
\[
     1_{\{\mathrm{failure}\}}e^{-Q}\le e^{-a}e^{-Q/2}.
\]
The reference integral of the remaining density must still be
estimated. For example replacing \(e^{-Q}\) by \(e^{-Q/2}\) in
a full-space \(d\)-dimensional positive quadratic Gaussian integral multiplies
its mass by \(2^{d/2}\). This follows by scaling all \(d\)
coordinates by \(\sqrt2\) and is a volume cost to be recorded.
For a site-gauge factor normalized by the truncated Haar number
$z(y)$ of \cref{rt:lem-haar-normalizer}, the corresponding ratio
after half of its energy is extracted is at most
$z_{\rm full}(y/2)/z(y)\le C$ for $y\ge2$. Here $z_{\rm full}$
uses cutoff one. The upper and lower $y^{-3/2}$ estimates in that
lemma prove this bound, so this normalization spends a bounded
logarithmic price at each such site.
These uses of positive energy require an allocation of disjoint
shares, including multiplicities when carriers overlap. The same
energy share cannot produce a gain in two factors of
\eqref{rt:eq-retained-product}.

\subsection{Counting the standing labels}
\label{rt:subsec-history-inputs}

We state the local quantitative inputs for the label estimate.
They concern the exact zero-source representation at the stated
stop. They are supplied by the regular and large-field steps on
the actual carriers and reference laws. No bound on the sum over
final supports is part of these inputs.

The cell side at level \(t\) is \(M\check R_t\) in its
lattice units. Write \(\vartheta_t=\min_{i\le t}\theta_i\).
In this partition estimate specialize the cell exponent of
\eqref{rt:eq-running-cell} to a fixed \(r\ge1\). Use the actual rounded scale
\[
      \check R_t=L^{\lceil \log_L((\log\vartheta_t)^r)\rceil}.
\]
The frame bound \(\theta_i-\theta_t\ge-c_*\) for \(i\le t\)
gives \(\theta_t-c_*\le\vartheta_t\le\theta_t\).
Comparability \(x_t\asymp\theta_t\) on the real tube
therefore gives \(1+\log x_t\le C\check R_t\) above a fixed
floor, while \(\check R_t\le L(\log\theta_t)^r\).
If \(\theta_t-\theta_{t+1}\le B\), the running minimum
drops by at most \(B\) in a step. Above a further fixed floor,
\((\log\vartheta_t/\log(\vartheta_t-2B))^r<L\).
Consequently the rounded scale decreases by at most one
power of \(L\) over two successive steps. The physical
partitions, which are aligned, either agree or coarsen by
\(L\) at each step, and two successive partitions cannot
both agree with their predecessors.

The final cells form a finite graph by the adjacency used in
the operation. Its maximum degree \(D_0\) is fixed, and the
number of cells is at most \(N_n\). A label has a unique
collection of connected standing supports in that graph.
For a fixed collection, its remaining visible information is
recorded separately under those supports. Internal records
already summed in a retained mass need only be tracked insofar
as the label distinguishes them. If labels distinguish their
visible traces, the mass in \eqref{rt:eq-weighted-retained-mass}
is split according to those traces; summing the disjoint parts
recovers the original mass. The record for a support may be
bounded by a larger set of possibilities, with repetitions only
if their counts are charged explicitly.

For each support \(C\), root its record at a fixed final cell
of \(C\). The active addresses are level-cell pairs, including
the final cells of \(C\). They form a connected set by the
parent and merger rules, with any necessary bounded collars
included in the charge. At successive levels the partition
either stays the same or coarsens by \(L\), so a cell has at
most \(L^4\) children. A neighbor at the same level has bounded
range in cell units. The address graph has fixed degree
\(D\ge2\), also on a periodic quotient. For an active set
\(\mathcal A_C\) define its charged volume by
\[
 v_C=\sum_{(t,\mathfrak c)\in\mathcal A_C}
       (M\check R_t)^4(1+\log x_t).
\]
This dominates both the number of its active addresses and
\(|C|\), the number of final cells in \(C\). A fixed
multiplicity of charges changes the cost constant. For a
fixed active set the count of the other
record choices, excluding normalized frame choices, is at
most \(\exp(bv_C)\) with a fixed \(b\). In the actual local
code this follows when a bounded alphabet records a bounded
number of tests per site and bond, and choices of stage or
failure type at level \(t\) have at most a fixed power of
\(1+\check R_t\) possibilities per active cell.
Indeed \(\log(1+\check R_t)\le C(1+\log x_t)\) for the
chosen polynomial sizes in \(\log x_t\). Multiplying the
local counts gives the stated bound with the own-scale volume.
Frame counts are integrated with their normalized laws, which
have total mass at most one on any restriction. Every
additional discrete choice must be included in this local code.

The second input is the local action and cost estimate. Separate
the exponential remaining after removing \(v_{n,\mathcal S}\)
into the assigned energy factors, the retained integrands and
the remaining regular, remainder and boundary exponent. Its
positive part, together with the unused reference and vacuum
costs, is bounded by
\begin{equation}\label{rt:eq-history-cost}
      A_0(1+\log x_n)N_n + A_1\sum_C v_C.
\end{equation}
The first term is the contribution of the regular region in
final units. The second is an own-scale estimate for terms
meeting the recorded large-field regions. In particular the
missing vacuum contributions are charged at their own levels:
each omitted or reassigned numerical contribution has a charged
carrier of bounded multiplicity and costs at most a constant
times \(1+\log x_t\) per site at its level \(t\).
The same allocation applies to normalization costs of the
reference integrals. Any part of the cost inside a retained
operation is paid using \eqref{rt:eq-weighted-retained-mass}
at a reserve weight chosen before the inverse-coupling floor.
When a record sees only the visible part of a drop-out, the
remaining volume it sees is charged to that same operation
at the raised weight. Its already integrated creations are
not separate gains. Charges paid in this way can be removed
from \eqref{rt:eq-history-cost}.

After these choices, the actual weighted retained masses and
certified failed tests give a product of factors \(e^{-G_C}\).
For each fixed exterior record the analytic input consists of
the local failure gains and mass bounds with the supremum
in \eqref{rt:eq-weighted-retained-mass}. In a component
they supply gains \(g(e)\) on its paying events, so
\(G_C=\sum_e g(e)\). Here \(e\) ranges over creations outside
outermost retained carriers and the distinct retained marks
with one certified, unused gain assigned to their visible successor. A creation outside
those carriers supplies \(g(e)\ge c_1 p_0(g_t^{\rm ref})^2\) in its
own level \(t\), or the reduced \(c_1p_0(g_t^{\rm ref})\) estimate
after any preceding allocations. A mark supplies \(g(e)\ge c_1p_0(g_t^{\rm ref})\) either from a
reserved part of its weighted mass or from a separately allocated
last-test share. For a native drop-out, its inner weighted mass
excludes the failed last test: that test supplies the gain for
the renewed wait, while the inner mass pays its recorded loads.
No gain is inferred merely from the presence of a mark, and each
energy share is used once.
We use \(g_t^{\rm ref}=\theta_t^{-1/2}\) and
\(p_0(g_t^{\rm ref})=A(\log\theta_t)^p\) for fixed \(A,p>0\).
Every nonempty component has at least one paying event.

The geometric inputs for the genealogy charge are the following
features of the actual hull and waiting rules. Each active
component has a skeleton of occupied cells connected by
bounded-range steps in the cell graph, and each charged cell
is within a fixed cell distance of the skeleton. At the next
level the occupied cells come from the parent images of
predecessor skeletons and at most a bounded number of new
occupied cells per paying event. Here a retained mark uses
the visible skeleton after its own hidden records have been
integrated. The seeds of an event enter at its own level or
at a later level, as for a bare failure recorded at the first
step. A cut leaf that needs a new gain uses an assigned
last-test share outside the inner mass.

The fixed-distance property is compatible with a collar
at every step. If an old collar has radius \(d\) in its
cell units and the partition coarsens by \(q\), its
new radius is at most \(d/q+A\) for a fixed \(A\),
including the parent-boundary ambiguity and the new collar.
Over two levels the older radius contracts by \(L^{-1}\).
Its accumulated contribution is therefore bounded by
\(2A/(1-L^{-1})\). Boxing only components of bounded diameter
adds a fixed amount to \(A\). The unique-address parent
map is applied to the occupied cells, while the closed
cell carriers and their collars enter the volume estimate.

A stretch with one predecessor and no new
paying event evolves by the quiet rule. For a fixed
skeleton-size threshold, once the skeleton has at most
that many cells, at most \(C\check R_t\) quiet levels may
occur before another event, a join, the end of the genealogy,
or a healing,
if the stretch started at level \(t\).
The constant may depend on the fixed threshold. This is
the local waiting bound, with the size at the start of the
stretch. We use it only on the non-erased
genealogy. Each unerased predecessor has a unique successor
containing its whole hull; components can merge but cannot split.
Consequently the ancestry of one standing component is a rooted
tree, directed toward that component. Every leaf has a paying
event. The size
\(\check R_t\) is nonincreasing, and the partition has no
two consecutive steps without coarsening.

The waiting bound uses an embedded light component and its
required collar. If that light geometry has no embedded
chart, its bounded cell diameter and bounded collars imply
that some torus period in cell units is bounded. This
period is divided by at least \(L\) every two further
levels. Since it must stay an integer at least one, only
a fixed number of further levels can be admitted. This is
also covered by the waiting bound, by increasing its
constant.

\begin{lemma}[Own-scale charge of a standing genealogy]
\label{rt:lem-genealogy-charge}
Under the skeleton and waiting conditions above, the record
volume is bounded by
\begin{equation}\label{rt:eq-genealogy-charge}
        v_C\le C_{\mathrm{geo}}\sum_e M^4\check R_{t(e)}^6 .
\end{equation}
The sum uses the same paying events as the mass estimate
and has no factor depending on the total depth.
\end{lemma}
\begin{proof}
Over \(s\) steps the parent partition coarsens by a factor
\(Q\ge L^{\lfloor s/2\rfloor}\) in cell units. For a connected
skeleton of \(m\) cells, walk twice along a spanning tree.
The walk has total length at most \(C m\) in the old cell
units, with every step of bounded length. Divide it into at
most \(C(1+m/Q)\) consecutive portions of length at most
\(CQ\). Each portion lies in a box of side \(O(Q)\), hence
meets only \(O(1)\) parent cells. On a torus lift each portion
before counting its parent cells; projection can only
decrease the number. Thus the skeleton has at most
\(C_1(1+m/Q)\) parent cells. Fix \(s\) so that \(Q\ge4C_1\).
For \(m\ge4C_1\) the count is at most \(m/2\); below
this fixed size it is bounded by a constant.

Write \(o_j\) for the total number of occupied cells in
the level-\(j\) members of the genealogy, \(l_j\) for the
number of its members below the fixed size, and \(a_j\)
for the number of events whose seeds enter at \(j\). The parent rule,
allowing a new occupied cell per event with fixed multiplicity,
gives
\[
   o_j\le \tfrac12 o_{j-s}+C l_{j-s}
                      +C\sum_{j-s<t\le j}a_t .
\]
Use zero counts before the first level. For
\(r_j=M^4\check R_j^5\), multiply by \(r_j\) and sum.
Because \(r_j\) is nonincreasing, moving the first sum to the
left gives
\[
       \sum_j r_j o_j
          \le C\sum_j r_j l_j+C\sum_t r_t a_t .
\]
If an event pays at an earlier level than its insertion,
monotonicity bounds its weight by the weight at its paying
level. Bounded skeleton neighborhoods and
\(1+\log x_j\le C\check R_j\) bound \(v_C\) by a constant
times the left side.

Divide the genealogy tree into stretches by beginning a new
stretch at a paying event, just after a keeping mark, or at
a join. Along a stretch without
an event the parent images cannot increase the number of
occupied cells. Its light members therefore form a last
segment, of length at most \(C\check R_t\) for a start at
\(t\), by the quiet waiting bound. Its cost in
\(\sum_j r_j l_j\) is at most \(C r_t\check R_t\).
Charge an event start or a successor of a mark to its event
or that preceding mark, respectively. To charge
join starts, designate a first predecessor at each vertex
with predecessors. At a join take a non-first predecessor,
then follow first predecessors down to a leaf. This assigns
distinct leaves to distinct joins. Indeed on the path up
from the chosen leaf the first ascent from a non-first
predecessor occurs at the specified join. A leaf is no later
than the join and has a paying event. Monotonicity of
\(r_t\check R_t\) allows that earlier event to pay.
Each event is used at most a fixed number of times.
This bounds the right side of the last display by the right
side of \eqref{rt:eq-genealogy-charge}.
\end{proof}

Different histories with one visible label can have different
exterior creation sets. Extracting their largest remaining factor
requires a common reference after all the preceding changes of
variables and density factors. More precisely, suppose the hidden
record $h$ is integrated by a nonnegative kernel $K(dh\mid z)$
with $K(1\mid z)\le1$. If $H(h,z)\le M(z)$ on its support,
then $\int H(h,z)K(dh\mid z)\le M(z)$. The bound $M(z)$
remains inside the $z$ integral. To select a single worst record
for the label, first extract a common nonnegative integrable
envelope in $z$; the remaining outside factor, including its
record-dependent numerical costs, must depend only on $h$.
Its maximum is then a deterministic factor. The chosen record
and visible trace must remain in the stated code count. The original stochastic operation alone
does not provide this domination after a density is refactored.
An alternative is to keep the actual disjoint record pieces and
include their reference masses or multiplicities in the count
before summing. One of these two warrants is part of the actual
record input. Without it, $m$ pieces of mass $e^{-G}$ have total
mass $m e^{-G}$, which a maximum alone does not bound.

\begin{lemma}[Rooted record count]\label{rt:lem-history-code}
For each final support \(C\), choose a fixed root in \(C\).
Suppose the record inputs just stated hold. There is a fixed
\(c_{\rm e}\) for which the sum of \(e^{-c_{\rm e}v_C}\) over
its possible outer records, with normalized frame laws, is
at most one.
\end{lemma}
\begin{proof}
Order the neighbors of each vertex of the address graph.
Associate to each connected active set of size \(m\) containing
the root a deterministically chosen rooted spanning tree, then
record the depth-first walk that traverses each tree edge twice.
The visited set determines the active set, and the walk has
length \(2(m-1)\). Hence there are at most \(D^{2(m-1)}\)
choices of active set. The other choices for an active set
use at most \(e^{b v_C}\) codes. Thus the sum
is at most
\[
   \sum_{m\ge1}D^{2(m-1)}e^{-(c_{\rm e}-b)m}\le 1
\]
if \(c_{\rm e}-b\ge 2\log D+\log2\).
The root level is fixed. Positions at earlier levels are
part of the same walk, and do not cost a factor equal to
the depth or the fine volume.
\end{proof}

\begin{proposition}[Upper bound for the history sum]
\label{rt:prop-history-upper}
Suppose the zero-source representation satisfies
\eqref{rt:eq-history-decomposition}, the integral and gain inputs
of \cref{rt:subsec-retained-integrals}, and the record, cost and
genealogy inputs of \cref{rt:subsec-history-inputs}. Assume
\(\check R_t\le C(1+\log\theta_t)^r\), \(p>6r\), and
\(\theta_t\ge\theta_n-c_*\) for \(t\le n\), with a fixed
\(c_*\). Then, with the inverse-coupling floor sufficiently
large, uniformly in \(n\), cutoff and admissible volume,
\begin{equation}\label{rt:eq-partition-upper}
    \widehat Z_n=e^{-E_{(n)}}Z_{K',\ell,\beta}
          \le \exp\{C(1+\log x_n)N_n\}.
\end{equation}
The same bound holds for the sum of absolute retained integrals
with any inserted component weights already covered by the
specified reserve.
\end{proposition}
\begin{proof}
Insert and remove the numerical factor \(e^{v_{n,\mathcal S}}\)
in the common normalization of each term of
\eqref{rt:eq-history-decomposition}. Bound the residual
exponent by \eqref{rt:eq-history-cost}. The retained and
failure factors are handled by \cref{rt:lem-component-majorant}
with the reference product of the actual record, including the
reserve for any inserted weight. The finite sums over disjoint
pieces of an internal record can be taken inside the resulting
majorants. We obtain an upper bound by a sum of products, with
factor \(e^{A_1v_C-G_C}\) for each outer record of \(C\), and
the same factor \(e^{A_0(1+\log x_n)N_n}\) for all labels.

Choose \(\kappa>2\log D_0\), using \(D_0\ge2\) without loss
of generality. Since \(p>6r\), the gain of each paying event
exceeds \(2(A_1+\kappa+c_{\rm e})C_{\rm geo}
M^4\check R_{t(e)}^6\) when the common floor is large
enough. The floor can also ensure
\(\log\theta_t\ge\tfrac12\log\theta_n\) for \(t\le n\).
Thus one half of \(G_C\) pays \((A_1+\kappa+c_{\rm e})v_C\)
by \eqref{rt:eq-genealogy-charge}. The other half leaves
\(s_n=c_{\rm gain}p_0(g_n^{\rm ref})\) for a fixed \(c_{\rm gain}>0\), because there
is at least one paying event. It follows that
\[
       e^{A_1v_C-G_C}
        \le e^{-s_n-\kappa|C|}e^{-c_{\rm e}v_C}.
\]
The sum over records on a fixed \(C\) is at most
\(a(C)=e^{-s_n-\kappa|C|}\) by \cref{rt:lem-history-code}.
This step takes place before the sum over placements of the
final supports.

Apply the same walk count in the final graph, now with
degree \(D_0\). The sum of \(a(C)\) over all nonempty
connected supports is at most
\[
   N_n e^{-s_n}\sum_{m\ge1}D_0^{2(m-1)}e^{-\kappa m}
       \le C_0 N_n .
\]
Finally, discard compatibility between distinct final supports
in the sum over their collections. Every collection occurs
at most once once its records are fixed. Its weight is
nonnegative, so the sum is at most
\(\prod_C(1+a(C))\le\exp(\sum_C a(C))\).
Combining the bounds proves \eqref{rt:eq-partition-upper}.
\end{proof}

The reference to actual component integrals in this argument
cannot be replaced by a pointwise bound on one polymer gas.
If a representation is \(Z=\int q(y)\Xi_y\,d\pi(y)\), where
\(q\ge0\) and a real logarithm satisfies
\(|\log\Xi_y|\le cN_n\), then
\[
   e^{-cN_n}\int q\,d\pi\le Z\le e^{cN_n}\int q\,d\pi.
\]
For an empty gas and \(q\) the indicator of a Haar set of
measure \(e^{-N_n^2}\), the logarithm vanishes and \(Z=e^{-N_n^2}\).
The lower bound below uses an explicit event with a controlled
Haar measure.

\subsection{The upper bound with a specified label load}

The actual read ownership gives a more precise estimate when the
remaining unmarked regular energy has an extensive bound in final
units. We state this version with the scalar load in its own level
units. On a standing label let \(\Omega_j\) be the determining
regions and let \(\Lambda_j^{\mathrm a}\) be the small-field regions
of the action sequence. Put \(Z_j^{\mathrm a}=(\Lambda_j^{\mathrm a})^c\)
and let \(Z_j^+\) be its one-cell enlargement. Write
\(\#_j(D)\) for the number of level-\(j\) sites in \(D\).
Set \(Y_0^{\mathrm d}=T_\ell\setminus\Omega_1\) and
\(Y_m^{\mathrm d}=\Omega_m\setminus\Omega_{m+1}\) for \(1\le m\le n\),
with \(\Omega_{n+1}=\varnothing\). Thus \(Y_0^{\mathrm d}\) is the
full outside layer in bare units. For \(c\ge0\), define
\begin{equation}\label{rt:eq-label-load}
 \mathscr L_c(\mathcal S)
   =c\sum_{m=0}^{n}\#_m(Y_m^{\mathrm d})
      +\sum_{j=1}^{n}
       \{C_{\log}\log\theta_{j-1}+C_{\rm loc}+c\}\#_j(Z_j^+).
\end{equation}
At \(n=0\) use \(Y_0^{\mathrm d}=T_\ell\). For each fixed $c$, the two price constants in
\eqref{rt:eq-label-load} are nonnegative. Fix them and the raised
retained-integral reserve before the inverse-coupling floor. Use \(\sigma_{\mathcal S,c}\) for the
literal absolute integral of this label after stripping its
own factor \(e^{v_{n,\mathcal S}}\) and inserting the factors
\(e^{c\,v(\omega)}\) at the outermost retained operations.

\begin{proposition}[The loaded history estimate]
\label{rt:prop-loaded-upper}
Use the actual product coordinates, reserve, and record inputs
of \cref{rt:prop-history-upper}. Suppose the entire remaining
final-site price after stripping the label vacuum is bounded by
\(C_{\rm sf}N_n\). This includes the unmarked energy and every
regular, remainder, boundary or reference factor not already paid
by the retained masses or assigned own-scale defect prices.
The actual remaining Wilson factor uses the ordered profiles of
\eqref{rt:eq-positive-action-profile}, and its allocated shares
leave nonnegative plaquette coefficients.
The determining regions
are nested with \(\Omega_{m+1}\) aligned at level \(m+1\),
and \(\Lambda_j^{\mathrm a}\subset\Omega_j\).
The frame has \(1<\theta_j\) and
\(\theta_{j-1}-\theta_j\le B\) with a fixed \(B\).
Use the disjoint final-cell partition from the standing geometry,
in which every cell contains at least \(M^4\) level-\(n\) sites.
Assume the actual local read ownership. Every earlier layer
address and enlarged defect address lies in the collar of its
assigned action cell, or is assigned to the latest retained
operation that remade that address. These assignments have
bounded multiplicity in the original level of each address.
For a remade portion of volume \(V_{\rm in}\), measured
with its own logarithmic site weights, every retained record
allowed by the visible label has \(v(\omega)\ge V_{\rm in}\).
Assume the local mass bounds hold at the correspondingly raised
reserve, chosen before the inverse-coupling floor. Then
\begin{equation}\label{rt:eq-loaded-label-sum}
 \sum_{\mathcal S} e^{\mathscr L_c(\mathcal S)}
                   \sigma_{\mathcal S,c}(1)
       \le \exp\{(c+C_{\rm sf}+M^{-4})N_n\}.
\end{equation}
If the price \(\mathscr L_0\) bounds the positive omitted
vacuum number \(v_{n,\mathcal S}-E_{(n)}\), this also bounds
\(\widehat Z_n\) by the right side for \(c=0\).
\end{proposition}
\begin{proof}
The last layer contributes at most \(cN_n\). For \(m<n\),
\[
 \#_m(Y_m^{\mathrm d})\le\#_m(\Omega_{m+1}^c)
      =L^4\#_{m+1}(\Omega_{m+1}^c)
      \le L^4\#_{m+1}(Z_{m+1}^{\mathrm a}).
\]
The equality counts the level-\(m\) sites of each aligned
level-\((m+1)\) block. The first inequality at \(m=0\)
includes every bare site of the outside layer. Assign
the enlarged \(Z_j^+\) cells to neighboring original
cells of \(Z_j^{\mathrm a}\), with bounded multiplicity.
The one-step frame bound gives
\(\log\theta_{j-1}\le C(1+\log\theta_j)\).
It follows that the earlier portions of \(\mathscr L_c\)
are bounded by own-scale prices, with constants affine in
\(c\). Apply the read ownership to retain the assigned
addresses when an operation remakes a determining layer.
An ordinary address is paid from the genealogy's event
reserve. A remade address is charged to the operation that
owns it, in the original level units, even if several later
terms read that operation.

For the latter cost first restrict to the compatible
inner records \(\mathcal B\) of that mark. If its assigned
price is at most \(A_c V_{\rm in}\), choose \(c'\ge c+A_c\).
On these records
\[
    e^{A_cV_{\rm in}}e^{c v(\omega)}
                          \le e^{c' v(\omega)} .
\]
Consequently the price times \(\mathsf M_c(i,X;\mathcal B)\)
is bounded by \(\mathsf M_{c'}(i,X;\mathcal B)\), with the
same restriction. The unrestricted raised mass is an upper
bound for this restricted mass. Keeping the restriction when
labels separate visible records preserves their disjoint sums.
The older failures already integrated in the mark are not
charged or used as a gain a second time. An unfinished live
lineage has no completed mark mass; its actual failures and waiting
costs are charged by the preceding genealogy estimate.

Apply the exterior-supremum estimate and the event prepayment
of \cref{rt:prop-history-upper} with these fixed, increased
local costs. The surviving regular factor is at most
\(e^{(c+C_{\rm sf})N_n}\). By
\eqref{rt:eq-positive-action-profile} and the share allocation,
the remaining real Wilson factor is at most one. The remaining Haar probability
measure has mass at most one on the cut. Each final connected
component \(C\), after summation of its records, keeps the
bound \(e^{-s_n-\kappa |C|}\). Choose the fixed \(\kappa\)
so that \(\sum_{m\ge1}D_0^{2(m-1)}e^{-\kappa m}\le1\).
There are at most \(N_n/M^4\) possible root cells, since a
final cell contains at least \(M^4\) sites. The sum of these
component weights is at most \(N_n/M^4\). The product over
distinct compatible components is bounded by its exponential,
as in the proof of \cref{rt:prop-history-upper}. This gives
\eqref{rt:eq-loaded-label-sum}. At $c=0$ the
absolute integral controls the vacuum-stripped label, so
the final assertion follows by reinserting the common
vacuum difference.
\end{proof}

\subsection{An event in the empty standing label}

For the empty label the determining sequence is constant and
only the final bonds \(V\) remain. Use the actual small-field
test \(\chi_n\) and background \(U_n(V)\) of this label. The
supplied density formula is
\begin{equation}\label{rt:eq-empty-label-density}
 e^{-E_{(n)}}d\tau_{\emptyset}(V)
   =\chi_n(V)\exp\{-x_n A(U_n(V))+F_n(U_n(V))\}
                      \prod_{b}dV_b .
\end{equation}
The equality of the numerical factor here with \(E_{(n)}\)
uses the same localization and normalizers as
\eqref{rt:eq-vacuum-recursion}. On this label the exponent
\(F_n\) contains the vacuum-subtracted regular and remainder
terms, including those produced by healing. There is no
standing boundary term.

We record the needed regularity inputs for this formula.
For age \(h=n-j\), the bound on the pointed, vacuum- and
classically subtracted regular family is, per level-\(j\)
base site,
\[
  C_E\{(1+h)L^{-b_{\rm T}h}+L^{-h}\}L^{-4h},
  \qquad b_{\rm T}>0 .
\]
It is evaluated on the constant determining sequence, where
the partition functions used for pointing sum to one. The
corresponding bound on the vacuum-subtracted remainder is
\(C_R(1+h)L^{-4h}\theta_j^{-q}\), with \(q>2\).
These are rooted bounds on the actual test support. They include
the localization tails on wrapping domains and any terms made
by the large-field step. Assume
\(\theta_j\ge\theta_n+\lambda(n-j)-c_*\) for \(\lambda>0\)
and \(\theta_n\ge1+c_*\). There are \(L^{4h}N_n\) base sites.
Summing first at age \(h\) removes \(L^{-4h}\); summing over
ages uses the two convergent series
\[
   \sum_{h\ge0}\{(1+h)L^{-b_{\rm T}h}+L^{-h}\},
   \qquad \sum_{h\ge0}(1+h)(1+\lambda h)^{-q}.
\]
Consequently \(|F_n|\le C_I N_n\) on the stated support.

\rtneedspace{25\baselineskip}
\begin{lemma}[Admitting a product of link balls]
\label{rt:lem-small-field-event}
Suppose the actual terminal test uses the separated
determining towers of the empty standing label, with all
tested plaquettes in the interior test window. The actual
determining datum on each tower, including the constant global
sequence, is the centered face lift of the terminal bonds and
its prescribed iterated averages, as described in the proof.
This is an identity for the datum used by the operation. Assume the
linear background regularity estimate on these towers and
on the constant global sequence. Namely, a determining datum
of plaquette radius \(d\) below a fixed ceiling gives
plaquette defects at most \(B_3dL^{-2n}\) on the tested
plaquettes and on the bare plaquettes of the global
background, respectively. Suppose $\chi_n$ is exactly the product
of these terminal plaquette tests on the admitted stock; any
other terminal condition has already been verified to be one
there. Each plaquette test is one below \(\varepsilon_n L^{-2n}\), with
\(\inf_n\varepsilon_n\sqrt{\theta_n}>0\) and
\(x_n\asymp\theta_n\ge1\) uniformly. Then for some fixed
\(c_{\rm sf}\in(0,1]\), on
\[
   \mathcal B_n=\{V:\|V_b-I\|_{\mathrm{op}}
                         <c_{\rm sf} x_n^{-1/2}\text{ for every }b\}
\]
one has \(\chi_n=1\) and
\begin{equation}\label{rt:eq-floor-action}
   x_n A(U_n(V))\le 48B_3^2 c_{\rm sf}^2 N_n .
\end{equation}
At level zero the action estimate holds with \(B_3=1\)
without a background construction.
\end{lemma}
\begin{proof}
Put \(s=c_{\rm sf}x_n^{-1/2}\). The product of the four unitary
matrices around a final plaquette differs from the identity
in operator norm by at most \(4s\), by telescoping and
unitarity. To compute the radius of the actual determining
datum, lift \(V\) to finer links as follows. A link staying
inside a top block is the identity. A link crossing a top
block face is assigned its corresponding oriented top link.
These prescriptions use the periodic centered blocks, or
their common lift in a test window, and commute with reversal.

An elementary fine plaquette crossing faces in no coordinate
has holonomy \(I\). If it crosses faces in only one coordinate,
the two crossing factors are inverses and the other two are
identities. If it crosses in both coordinates, the four
crossing factors form the corresponding top plaquette
holonomy with its induced orientation. Thus its radius is
at most \(4s\). The one-step averaging of this link field
is its lift at the coarser resolution. Indeed, at an
intermediate scale the two endpoint contours are inside
their own top blocks and contribute identities. Each joining
segment on a coarse bond crosses either no top face or
the same single face in the same orientation, so every
transported candidate in the group mean has the same value.
Centered nested blocks ensure this also for projected paths
at a periodic face, counting the oriented path occurrences.
Iteration gives the same radius on the averaged rows of
each terminal tower. This is the radius of the actual datum
when it is formed by the prescribed lift and averages.

The stated linear estimates now give defect at most
\(4B_3sL^{-2n}\) on the tested stock and the global bare
background. Choose \(c_{\rm sf}\) small enough for \(4s\) to be
below the datum ceiling and \(4B_3s<\varepsilon_n\).
The choices are uniform by the two comparisons in the
statement. The test is then one. Earlier fluctuation tests
are already incorporated in \eqref{rt:eq-empty-label-density}.
For \(U\in\operatorname{SU}(2)\) with class angle \(\varphi\),
\[
      1-\tfrac12\operatorname{Tr}U
         =1-\cos\varphi=\tfrac12\|U-I\|_{\mathrm{op}}^2 .
\]
There are \(6N_0\) plaquettes in the bare Wilson sum and
\(N_0=L^{4n}N_n\). The two powers of \(L^{-2n}\) cancel
this site conversion and give
\(A(U_n(V))\le48B_3^2s^2N_n\).
At \(n=0\) use the direct plaquette telescoping with the
bare field \(V\) in the same calculation.
\end{proof}

\rtneedspace{23\baselineskip}
\begin{theorem}[Common-normalized partition bounds]
\label{rt:thm-partition}
Under the history hypotheses of \cref{rt:prop-history-upper}
and the empty-label density, rooted estimates and regularity
hypotheses of \cref{rt:lem-small-field-event}, one has
\begin{equation}\label{rt:eq-partition-two-sided}
   e^{-(C+6\log x_n)N_n}
     \ \le\ e^{-E_{(n)}} Z_{K',\ell,\beta}
     \ \le\ e^{C(1+\log x_n)N_n}.
\end{equation}
At the terminal first passage \(n=K\) with
\(h=g^{-2}\le x_K\le h+B_{\rm s}\) and \(K'=K\), this gives
\[
       e^{-C_g\ell^4}
          \le e^{-E_{(K)}}Z_{K',\ell,\beta}
          \le e^{C_g\ell^4}.
\]
The constants are independent of the bare cutoff and compatible
volume. At earlier stops \eqref{rt:eq-partition-two-sided}
uses \(N_n\) and requires the same local inputs at that stop.
\end{theorem}
\begin{proof}
For \(0<s\le1\), an \(\operatorname{SU}(2)\) matrix in
the class-angle range \(0\le\varphi<s\) belongs to the
operator-norm ball of radius \(s\), since
\(\|U-I\|_{\mathrm{op}}=2\sin(\varphi/2)\).
The normalized Haar angle density is
\((2/\pi)\sin^2\varphi\,d\varphi\). The inequality
\(\sin\varphi\ge 2\varphi/\pi\) for \(0\le\varphi\le1\)
therefore gives ball mass at least \(8s^3/(3\pi^3)\).
There are \(4N_n\) independently integrated positive links.

By positivity, keep the empty label in
\eqref{rt:eq-history-decomposition}, then restrict its
integral to \(\mathcal B_n\). Use \(|F_n|\le C_I N_n\)
and \eqref{rt:eq-floor-action}. With \(c_H=8/(3\pi^3)\)
this gives the lower bound
\[
  e^{-E_{(n)}}Z_{K',\ell,\beta}
     \ge e^{-(C_I+48B_3^2c_{\rm sf}^2)N_n}
            (c_Hc_{\rm sf}^3 x_n^{-3/2})^{4N_n}.
\]
The upper bound is \eqref{rt:eq-partition-upper}.
At the first passage with \(K'=K\), \(N_K=\ell^4\) and the inverse
coupling lies in a fixed interval depending on \(g\).
Absorb its endpoint into \(C_g\).
\end{proof}

\begin{thebibliography}{9}
\raggedright
\bibitem{EM} \emph{Iteration of regular interactions on prescribed coupling prefixes}, companion manuscript, 2026.
\bibitem{Bal88} T. Ba\l aban, \emph{Convergent renormalization expansions for lattice gauge theories}, Comm. Math. Phys. \textbf{119} (1988), 243--285. \href{https://doi.org/10.1007/BF01217741}{doi:10.1007/BF01217741}.
\bibitem{Bal89} T. Ba\l aban, \emph{Large field renormalization. II. Localization, exponentiation, and bounds for the R operation}, Comm. Math. Phys. \textbf{122} (1989), 355--392. \href{https://doi.org/10.1007/BF01238433}{doi:10.1007/BF01238433}.
\end{thebibliography}
\end{document}
